\documentclass[12pt,a4paper,UTF8]{ctexart}
\usepackage[a4paper,top=22mm,bottom=22mm,left=22mm,right=22mm,headheight=15pt]{geometry}
\usepackage{amsmath,amssymb,mathtools,bm}
\usepackage{graphicx,xcolor,booktabs,longtable,array,enumitem,fancyhdr,caption,float,microtype,setspace}
\usepackage[hidelinks,unicode]{hyperref}
\usepackage[most]{tcolorbox}
\definecolor{ContestBlue}{HTML}{17365D}
\hypersetup{pdftitle={第083套},pdfauthor={}}
\setstretch{1.24}
\setlength{\parindent}{2em}
\setlength{\parskip}{0.35em}
\setlength{\emergencystretch}{2em}
\setlist{itemsep=0.35em,topsep=0.4em}
\captionsetup{font=small,labelfont=bf,justification=centering}
\setcounter{secnumdepth}{0}
\pagestyle{fancy}
\fancyhf{}
\fancyhead[L]{\small 第083套}
\fancyfoot[C]{\thepage}
\renewcommand{\headrulewidth}{0.4pt}
\ctexset{
  section={format=\Large\bfseries\color{ContestBlue},beforeskip=1.2em,afterskip=0.7em},
  subsection={format=\large\bfseries\color{ContestBlue},beforeskip=1.0em,afterskip=0.5em},
  subsubsection={format=\normalsize\bfseries,beforeskip=0.8em,afterskip=0.35em}
}
\makeatletter
\newcommand{\tightlist}{\setlength{\itemsep}{0pt}\setlength{\parskip}{0pt}}
\newcommand{\maxwidth}{\ifdim\Gin@nat@width>\linewidth\linewidth\else\Gin@nat@width\fi}
\newcommand{\maxheight}{\ifdim\Gin@nat@height>0.72\textheight 0.72\textheight\else\Gin@nat@height\fi}
\newcommand{\pandocbounded}[1]{\begingroup\setkeys{Gin}{width=\maxwidth,height=\maxheight,keepaspectratio}#1\endgroup}
\makeatother
\begin{document}
\begin{titlepage}
  \centering
  \vspace*{58mm}
  {\Huge\bfseries\color{ContestBlue} 第083套\par}
  \vfill
  \vspace*{28mm}
\end{titlepage}
\section{第 1 题}\label{ux7b2c-1-ux9898}

在一颗重力加速度未知的星球上，有一根质量为 \(M\)、长度为 \(L\)
的细长空心圆柱。四个人均视为质点，质量均为
\(m\)，依次紧靠在圆柱内部，相邻两人的质心距离均为
\(s\)。最上方一人的质心到圆柱顶端的距离为 \(z\)。

实验测得整个系统的质心到圆柱顶端的距离为
\(x_{\rm cm}\)。在圆柱上选取两个小孔，它们到系统质心的距离分别为
\(R_1\)、\(R_2\)。使圆柱分别绕两个小孔作小角度摆动，测得周期分别为
\(T_1\)、\(T_2\)。

\begin{enumerate}
\def\labelenumi{\arabic{enumi}.}
\tightlist
\item
  给出一种测量系统质心位置的方法；
\item
  求该星球的重力加速度 \(g\)；
\item
  已知 \(M\)、\(m\) 时，求参数 \(z\)；
\item
  当 \(M\)、\(m\) 未知时，给出由实验量确定 \(z\)
  的中间过程，不要求写出最终展开式；
\item
  若可选择更多小孔并进行多次周期测量，请给出线性拟合公式。若拟合直线斜率为
  \(k\)，用 \(k\) 表示 \(g\)。
\end{enumerate}

\begin{figure}[H]
\centering
\includegraphics[width=0.472\linewidth,height=\textheight,keepaspectratio,alt={图 1-1　空心圆柱内四个质点的位置及主要几何量。}]{../figures/question-01.pdf}
\caption*{图 1-1　空心圆柱内四个质点的位置及主要几何量。}
\end{figure}

\subsection{一、题型定位}\label{ux4e00ux9898ux578bux5b9aux4f4d}

挖孔同时改变杆的质心和转动惯量。先由水平平衡条件定位质心，再把各悬点的复摆周期写成同一组质量矩，检查反解参数是否唯一。

\subsection{二、知识点}\label{ux4e8cux77e5ux8bc6ux70b9}

\begin{itemize}
\tightlist
\item
  离散质点与连续刚体组成系统的质心；
\item
  平行轴定理 \(I_P=I_C+\mathcal M R^2\)；
\item
  复摆的小角度振动；
\item
  两次测量消去未知转动惯量；
\item
  质量分数参数化与线性拟合。
\end{itemize}

\subsection{三、物理图像}\label{ux4e09ux7269ux7406ux56feux50cf}

把圆柱顶端作为竖直坐标原点，向下为正。圆柱本身的质心在
\(L/2\)，四个质点依次位于

\[
z,\qquad z+s,\qquad z+2s,\qquad z+3s.
\]

绕某个小孔摆动时，整个圆柱和四个质点是一个刚体。重力等效作用在系统质心
\(C\)，小孔 \(P\) 是固定转轴，\(PC=R\)。图 1-2
把质心测量与复摆测量分别画出。

\subsection{四、解题入口}\label{ux56dbux89e3ux9898ux5165ux53e3}

第一入口是把所有质量对圆柱顶端取一阶矩；第二入口是对小角度 \(\theta\)
写绕小孔的力矩方程。两次复摆实验中，系统总质量和质心转动惯量相同，只有
\(R\) 不同，因此可以相减消去转动惯量。

\subsection{五、解析}\label{ux4e94ux89e3ux6790}

\begin{figure}[H]
\centering
\includegraphics[width=0.827\linewidth,height=\textheight,keepaspectratio,alt={图 1-2　用平衡支点测量质心，并使系统绕小孔作复摆。}]{../figures/solution-01.pdf}
\caption*{图 1-2　用平衡支点测量质心，并使系统绕小孔作复摆。}
\end{figure}

图 1-2 左图 对应质心位置的直接测量；图 1-2 右图
对应复摆模型。后者中支点力 \(\bm F_P\) 通过转轴 \(P\)，对 \(P\)
的力矩为零；重力作用在系统质心 \(C\)，力臂由 \(PC=R\) 和偏角 \(\theta\)
决定。因此绕孔列转动方程时，只需保留重力矩。

\subsubsection{1.
质心位置的测量}\label{ux8d28ux5fc3ux4f4dux7f6eux7684ux6d4bux91cf}

一种直接方法是把圆柱水平放在刀口或细棱上，缓慢移动支承位置，直到系统既不顺时针转动也不逆时针转动。此时支点通过系统质心，测量支点到圆柱顶端的距离即可得到
\(x_{\rm cm}\)。

也可以把系统从两个不同位置悬挂，分别画出铅垂线；两条铅垂线的交点就是质心。对于本题的一维细长模型，水平平衡法更直接。

\subsubsection{2.
两个复摆周期与重力加速度}\label{ux4e24ux4e2aux590dux6446ux5468ux671fux4e0eux91cdux529bux52a0ux901fux5ea6}

记系统总质量为

\[
\mathcal M=M+4m,
\]

系统绕自身质心、垂直纸面的转动惯量为 \(I_C\)。小孔到质心的距离为 \(R\)
时，由平行轴定理，绕孔的转动惯量为

\[
I_P=I_C+\mathcal M R^2.
\]

系统偏离竖直方向小角度 \(\theta\) 时，重力对小孔的力矩为

\[
\tau=-\mathcal M gR\sin\theta.
\]

由图 1-2 右图，支点力作用线通过 \(P\)，故其力矩为零；重力作用线到 \(P\)
的垂直距离为 \(R\sin\theta\)，所以重力矩大小为
\(\mathcal M gR\sin\theta\)。负号表示力矩总是把系统拉回平衡位置。小角度下
\(\sin\theta\approx\theta\)，于是

\[
\left(I_C+\mathcal M R^2\right)\ddot\theta
=-\mathcal M gR\theta.
\]

整理成简谐振动标准形式：

\[
\ddot\theta+\omega^2\theta=0,
\qquad
\omega^2=\frac{\mathcal M gR}{I_C+\mathcal M R^2}.
\]

故周期为

\[
T=2\pi\sqrt{\frac{I_C+\mathcal M R^2}{\mathcal M gR}}.
\]

为减少式子长度，定义单位总质量的质心转动惯量

\[
h_I=\frac{I_C}{\mathcal M}.
\]

则

\[
\frac{gT^2}{4\pi^2}=R+\frac{h_I}{R}.
\]

分别代入两次实验：

\[
\frac{gT_1^2}{4\pi^2}=R_1+\frac{h_I}{R_1},
\qquad
\frac{gT_2^2}{4\pi^2}=R_2+\frac{h_I}{R_2}.
\]

两式分别乘以 \(R_1\)、\(R_2\)：

\[
\frac{gR_1T_1^2}{4\pi^2}=R_1^2+h_I,
\qquad
\frac{gR_2T_2^2}{4\pi^2}=R_2^2+h_I.
\]

相减后 \(h_I\) 消去：

\[
\frac{g}{4\pi^2}\left(R_1T_1^2-R_2T_2^2\right)
=R_1^2-R_2^2.
\]

因此

\[
\boxed{
g=\frac{4\pi^2\left(R_1^2-R_2^2\right)}
{R_1T_1^2-R_2T_2^2}
}.
\]

把 \(g\) 代回任一周期方程，还可得到

\[
\boxed{
h_I=
\frac{R_1R_2\left(R_1T_2^2-R_2T_1^2\right)}
{R_1T_1^2-R_2T_2^2}
}.
\]

这里 \(h_I\) 的量纲是长度平方，正好符合``转动惯量除以质量''的量纲。

\subsubsection{\texorpdfstring{3. 已知 \(M\)、\(m\) 时求
\(z\)}{3. 已知 M、m 时求 z}}\label{ux5df2ux77e5-mm-ux65f6ux6c42-z}

对圆柱顶端取一阶矩。圆柱的质量矩为 \(M(L/2)\)，四个质点的质量矩之和为

\[
\begin{aligned}
&m z+m(z+s)+m(z+2s)+m(z+3s)\\
&=m(4z+6s).
\end{aligned}
\]

所以

\[
x_{\rm cm}=\frac{ML/2+m(4z+6s)}{M+4m}.
\]

两边乘以 \(M+4m\)：

\[
(M+4m)x_{\rm cm}=\frac{ML}{2}+4mz+6ms.
\]

移项并除以 \(4m\)：

\[
\boxed{
z=\frac{2(M+4m)x_{\rm cm}-ML}{8m}-\frac{3s}{2}
}.
\]

\subsubsection{\texorpdfstring{4. \(M\)、\(m\) 未知时消元确定
\(z\)}{4. M、m 未知时消元确定 z}}\label{mm-ux672aux77e5ux65f6ux6d88ux5143ux786eux5b9a-z}

只知道质心位置还不能区分圆柱与四个人各占多少质量，因此必须把复摆实验给出的
\(h_I\) 也利用起来。

定义四个人的平均位置

\[
\bar x_p=z+\frac{3s}{2},
\]

以及四个人占总质量的比例

\[
w=\frac{4m}{M+4m}.
\]

圆柱占总质量的比例就是 \(1-w\)。质心式可写成

\[
x_{\rm cm}=(1-w)\frac{L}{2}+w\bar x_p.
\]

令

\[
\delta=x_{\rm cm}-\frac{L}{2},
\qquad
d_p=\bar x_p-\frac{L}{2},
\]

便有

\[
\delta=wd_p,
\qquad
w=\frac{\delta}{d_p}.
\]

接下来计算 \(h_I=I_C/\mathcal M\)。圆柱自身绕中心的转动惯量贡献为

\[
(1-w)\frac{L^2}{12}.
\]

四个质点相对其平均位置的偏移依次为

\[
-\frac{3s}{2},\quad-\frac{s}{2},\quad
\frac{s}{2},\quad\frac{3s}{2}.
\]

偏移平方和为

\[
\frac{9s^2}{4}+\frac{s^2}{4}+\frac{s^2}{4}+\frac{9s^2}{4}=5s^2.
\]

因此四个人内部离散分布对 \(h_I\) 的贡献为

\[
\frac{m}{M+4m}\,5s^2=\frac{5ws^2}{4}.
\]

圆柱质心与四人平均质心之间还产生``两群质量中心分离''的贡献。两群质量的质量分数为
\(1-w\) 与 \(w\)，中心间距为 \(d_p\)，故该项为

\[
w(1-w)d_p^2.
\]

于是

\[
h_I=(1-w)\frac{L^2}{12}+\frac{5ws^2}{4}+w(1-w)d_p^2.
\]

再代入 \(w=\delta/d_p\)：

\[
h_I=\frac{L^2}{12}
+\frac{\delta}{d_p}\left(\frac{5s^2}{4}-\frac{L^2}{12}\right)
+\delta(d_p-\delta).
\]

令

\[
C=\frac{5s^2}{4}-\frac{L^2}{12},
\qquad
A=\frac{h_I-L^2/12+\delta^2}{\delta},
\]

就得到只含 \(d_p\) 的方程

\[
d_p+\frac{C}{d_p}=A,
\]

即

\[
d_p^2-Ad_p+C=0.
\]

从这个二次方程求出 \(d_p\) 后，用

\[
z=d_p+\frac{L}{2}-\frac{3s}{2}
\]

即可得到 \(z\)。若二次方程给出两个数学根，应以 \(0<w<1\)
且四个质点均位于圆柱内部作为物理筛选条件。

\subsubsection{5.
多次测量的线性拟合}\label{ux591aux6b21ux6d4bux91cfux7684ux7ebfux6027ux62dfux5408}

从复摆周期公式出发：

\[
T^2=\frac{4\pi^2}{g}\left(R+\frac{h_I}{R}\right).
\]

两边乘以 \(R\)：

\[
T^2R=\frac{4\pi^2}{g}\left(R^2+h_I\right).
\]

改写为

\[
R^2=\frac{g}{4\pi^2}(T^2R)-h_I.
\]

以

\[
X=T^2R,
\qquad
Y=R^2
\]

作图，拟合式为

\[
Y=kX+b.
\]

比较系数可知

\[
k=\frac{g}{4\pi^2},
\qquad
b=-h_I.
\]

所以

\[
\boxed{g=4\pi^2k}.
\]

\subsubsection{6.
复摆测量的自洽条件}\label{ux590dux6446ux6d4bux91cfux7684ux81eaux6d3dux6761ux4ef6}

由

\[
g=\frac{4\pi^2(R_1^2-R_2^2)}{R_1T_1^2-R_2T_2^2}
\]

可见，分子和分母必须同号，所得 \(g\) 才为正。若人为规定
\(R_1>R_2\)，实验数据应满足

\[
R_1T_1^2>R_2T_2^2.
\]

这不是额外假设，而是检验周期记录和孔距标记是否颠倒的办法。由实验反求的

\[
h_I=\frac{I_C}{M+4m}
\]

还必须满足
\(h_I>0\)。若得到负值，说明小角度条件、转轴摩擦或距离测量至少有一项不可信。

周期关于孔距并不是单调函数。由

\[
T^2=\frac{4\pi^2}{g}\left(R+\frac{h_I}{R}\right)
\]

可知，当 \(R=\sqrt{h_I}\) 时周期最小。若 \(R_1\)、\(R_2\)
选得过于接近，差分公式的分子、分母都很小，测量误差会被放大。因此实验上应选择相隔较明显、但又能保持稳定小摆的两个孔。

未知质量消元时，求出的质量分数

\[
w=\frac{\delta}{d_p}
\]

必须位于 \((0,1)\)。随后还要逐一检查

\[
0<z,
\qquad
z+3s<L,
\]

才能确认四个质点确实都在圆柱内部。这样，代数根、惯量根和几何构型才三者一致。

\textbf{反演的特殊情形与不唯一性}

正文在 \(\delta=x_{\rm cm}-L/2\ne0\) 时才可除以 \(\delta\)。若
\(\delta=0\) 且四人质量均为正，由 \(\delta=wd_p\) 得 \(d_p=0\)，直接有

\[z=L/2-3s/2.\]

随后 \(h_I=(1-w)L^2/12+w(5s^2/4)\)
可用来反推质量比例；若两个惯性尺度恰相等，周期不能再识别 \(w\)，但 \(z\)
仍已确定。对 \(\delta\ne0\)
的二次方程，物理边界筛选后也可能剩两个可行构型。若两根都满足 \(0<w<1\)
和内部位置限制，应报告两解并增加一项独立测量，不能任意选根。

图中水平平衡实验的重力竖直向下；转轴至质心的距离 \(R_i\)
与圆柱顶端至第一人的距离 \(z\) 不同，周期公式只使用前者。

\newpage

\section{第 2 题}\label{ux7b2c-2-ux9898}

一个半径为 \(R=20\,\mathrm{cm}\)
的光滑反射球可视为理想凸面镜。在球前放置一焦距为 \(f=20\,\mathrm{cm}\)
的薄凸透镜，透镜到球面顶点的距离为
\(d=15\,\mathrm{cm}\)。在透镜左侧、距透镜 \(u\) 处放置轴上点光源
\(S\)。光线依次经过凸透镜、凸面镜反射，再次经过凸透镜。

\begin{enumerate}
\def\labelenumi{\arabic{enumi}.}
\tightlist
\item
  为使最终成实像，求物距的下限 \(u_{\min}\)；
\item
  求物距的上限 \(u_{\max}\)；
\item
  把人的脸视为物体，并要求最终实像位于脸前 \(25\,\mathrm{cm}\)
  处。求脸到透镜的距离 \(u\)；
\item
  求第3问的总横向放大率，正值表示正立，负值表示倒立。
\end{enumerate}

\begin{figure}[H]
\centering
\includegraphics[width=0.83\linewidth,height=\textheight,keepaspectratio,alt={图 2-1　凸透镜、凸面反射球与轴上物点。}]{../figures/question-02.pdf}
\caption*{图 2-1　凸透镜、凸面反射球与轴上物点。}
\end{figure}

\subsection{一、题型定位}\label{ux4e00ux9898ux578bux5b9aux4f4d-1}

凸面镜位于透镜后，光先透射、再反射、再透射。给每次传播指定物像距符号，并用往返矩阵核对最终像的位置和放大率。

\subsection{二、知识点}\label{ux4e8cux77e5ux8bc6ux70b9-1}

\begin{itemize}
\tightlist
\item
  薄透镜成像公式；
\item
  球面镜焦距 \(f_m=-R/2\)；
\item
  实物、虚物、实像和虚像的符号；
\item
  多级光学系统的距离传递；
\item
  总横向放大率等于各级放大率之积。
\end{itemize}

\subsection{三、物理图像}\label{ux4e09ux7269ux7406ux56feux50cf-1}

设第一次通过透镜后的像距为 \(v_1\)；该像作为凸面镜的物，镜前物距为
\(u_2\)，反射像距为 \(v_2\)；反射光向左传播，再次进入透镜，其物距为
\(u_3\)，最终像距为 \(v_3\)。

对每个光学元件都采用``实物距、实像距为正''的局部符号约定。凸面镜焦距为

\[
f_m=-\frac{R}{2}=-10\,\mathrm{cm}.
\]

\subsection{四、解题入口}\label{ux56dbux89e3ux9898ux5165ux53e3-1}

先不要分类讨论每一级是实像还是虚像，而是直接把 \(v_1,u_2,v_2,u_3,v_3\)
全部表示成初始物距 \(u\) 的有理式。最终实像条件只需检查
\(v_3>0\)，这样能避免分支遗漏。

\subsection{五、解析}\label{ux4e94ux89e3ux6790-1}

\begin{figure}[H]
\centering
\includegraphics[width=0.9\linewidth,height=\textheight,keepaspectratio,alt={图 2-2　三次成像的位置传递及光线传播方向。}]{../figures/solution-02.pdf}
\caption*{图 2-2　三次成像的位置传递及光线传播方向。}
\end{figure}

图 2-2 按真实传播顺序标出了三个中间位置：第一次通过透镜得到
\(I_1\)，凸面镜反射得到 \(I_2\)，反射光向左第二次通过同一透镜得到最终像
\(I_3\)。图下方把距离传递写成
\(v_1\rightarrow u_2\rightarrow v_2\rightarrow u_3\rightarrow v_3\)；各物距、像距必须按对应光学元件的局部入射方向判断正负，不能沿用一条固定的全局正方向。

\subsubsection{1.
第一次通过凸透镜}\label{ux7b2cux4e00ux6b21ux901aux8fc7ux51f8ux900fux955c}

薄透镜公式为

\[
\frac{1}{u}+\frac{1}{v_1}=\frac{1}{f}.
\]

解得

\[
v_1=\frac{fu}{u-f}=\frac{20u}{u-20}.
\]

当 \(u<20\,\mathrm{cm}\)
时，\(v_1<0\)，第一次成虚像；这并不意味着整个系统不能最终成实像，因为后面还有反射和第二次折射。

\subsubsection{2. 凸面镜成像}\label{ux51f8ux9762ux955cux6210ux50cf}

透镜与球面顶点相距 \(d=15\,\mathrm{cm}\)，故凸面镜的物距为

\[
u_2=d-v_1.
\]

代入 \(v_1\)：

\[
\begin{aligned}
u_2
&=15-\frac{20u}{u-20}\\
&=\frac{15(u-20)-20u}{u-20}\\
&=-\frac{5(u+60)}{u-20}.
\end{aligned}
\]

凸面镜成像公式为

\[
\frac{1}{u_2}+\frac{1}{v_2}=-\frac{2}{R}=-\frac{1}{10}.
\]

因此

\[
v_2=\frac{-10u_2}{u_2+10}.
\]

先化简分母：

\[
\begin{aligned}
u_2+10
&=-\frac{5(u+60)}{u-20}+10\\
&=\frac{-5(u+60)+10(u-20)}{u-20}\\
&=\frac{5(u-100)}{u-20}.
\end{aligned}
\]

于是

\[
\begin{aligned}
v_2
&=-10\frac{-5(u+60)/(u-20)}{5(u-100)/(u-20)}\\
&=\frac{10(u+60)}{u-100}.
\end{aligned}
\]

\subsubsection{3.
第二次通过凸透镜}\label{ux7b2cux4e8cux6b21ux901aux8fc7ux51f8ux900fux955c}

反射光向左传播。相对于第二次入射方向，透镜物距为

\[
u_3=d-v_2.
\]

代入 \(v_2\)：

\[
\begin{aligned}
u_3
&=15-\frac{10(u+60)}{u-100}\\
&=\frac{15(u-100)-10(u+60)}{u-100}\\
&=\frac{5(u-420)}{u-100}.
\end{aligned}
\]

第二次透镜成像仍满足

\[
\frac{1}{u_3}+\frac{1}{v_3}=\frac{1}{20},
\]

所以

\[
v_3=\frac{20u_3}{u_3-20}.
\]

计算 \(u_3-20\)：

\[
\begin{aligned}
u_3-20
&=\frac{5(u-420)-20(u-100)}{u-100}\\
&=-\frac{5(3u+20)}{u-100}.
\end{aligned}
\]

因此最终像距可以写成非常简洁的形式：

\[
\boxed{
v_3=\frac{20(420-u)}{3u+20}
}.
\]

\subsubsection{4.
最终实像的物距范围}\label{ux6700ux7ec8ux5b9eux50cfux7684ux7269ux8dddux8303ux56f4}

实际物体位于透镜左侧，故 \(u>0\)。此时分母 \(3u+20\)
始终为正。要使最终像为实像，需要

\[
v_3>0,
\]

即

\[
420-u>0.
\]

所以严格的非退化范围是

\[
\boxed{0<u<420\,\mathrm{cm}}.
\]

因此物距的下、上界分别为

\[
\boxed{u_{\min}=0},
\qquad
\boxed{u_{\max}=420\,\mathrm{cm}}.
\]

这里 \(u_{\min}=0\)
是下确界，物体不能与理想薄透镜完全重合；\(u=420\,\mathrm{cm}\)
时最终像退化到透镜平面，也不属于通常所说的有限实像。

在 \(u=20\,\mathrm{cm}\) 或 \(u=100\,\mathrm{cm}\)
时，某一中间像位于无穷远，但总映射的极限仍然有限。不能因为中间公式出现分母为零，就把这些位置直接从最终成像范围中删除。

\subsubsection{\texorpdfstring{5. 像位于脸前
\(25\,\mathrm{cm}\)}{5. 像位于脸前 25\textbackslash,\textbackslash mathrm\{cm\}}}\label{ux50cfux4f4dux4e8eux8138ux524d-25mathrmcm}

脸在透镜左侧 \(u\) 处，最终实像也在透镜左侧，并位于脸与透镜之间，距脸
\(25\,\mathrm{cm}\)。因此最终像距为

\[
v_3=u-25.
\]

代入系统总成像式：

\[
\frac{20(420-u)}{3u+20}=u-25.
\]

交叉相乘：

\[
8400-20u=(u-25)(3u+20).
\]

展开右端：

\[
(u-25)(3u+20)=3u^2-55u-500.
\]

移到同一侧：

\[
3u^2-35u-8900=0.
\]

取正根：

\[
\begin{aligned}
u
&=\frac{35+\sqrt{35^2+4\times3\times8900}}{6}\\
&=\frac{35+\sqrt{108025}}{6}\\
&\approx60.612\,\mathrm{cm}.
\end{aligned}
\]

负根不符合 \(u>0\)，舍去。因此

\[
\boxed{u\approx60.61\,\mathrm{cm}}.
\]

\subsubsection{6.
总横向放大率}\label{ux603bux6a2aux5411ux653eux5927ux7387}

三次成像的横向放大率依次为

\[
\beta_1=-\frac{v_1}{u},
\qquad
\beta_2=-\frac{v_2}{u_2},
\qquad
\beta_3=-\frac{v_3}{u_3}.
\]

分别化简：

\[
\beta_1=-\frac{20}{u-20},
\]

\[
\beta_2
=-\frac{10(u+60)/(u-100)}{-5(u+60)/(u-20)}
=\frac{2(u-20)}{u-100},
\]

\[
\beta_3
=-\frac{20(420-u)/(3u+20)}{5(u-420)/(u-100)}
=\frac{4(u-100)}{3u+20}.
\]

三者相乘，大量中间因子消去：

\[
\begin{aligned}
\beta
&=\beta_1\beta_2\beta_3\\
&=-\frac{20}{u-20}
\cdot\frac{2(u-20)}{u-100}
\cdot\frac{4(u-100)}{3u+20}\\
&=-\frac{160}{3u+20}.
\end{aligned}
\]

代入 \(u\approx60.6119\,\mathrm{cm}\)：

\[
\boxed{\beta\approx-0.793}.
\]

负号表示最终像倒立，绝对值小于 \(1\) 表示像缩小。

\subsubsection{7.
中间成像位置的数值复核}\label{ux4e2dux95f4ux6210ux50cfux4f4dux7f6eux7684ux6570ux503cux590dux6838}

把 \(u=60.6119\,\mathrm{cm}\) 代回每一级成像，可得到

\[
v_1\approx29.8493\,\mathrm{cm}.
\]

因此对凸面镜而言

\[
u_2=15-v_1\approx-14.8493\,\mathrm{cm}.
\]

负物距表明到达球面的光线已经会聚，凸面镜面对的是虚物。继续计算：

\[
v_2\approx-30.6214\,\mathrm{cm}.
\]

于是第二次通过透镜时

\[
u_3=15-v_2\approx45.6214\,\mathrm{cm}>f.
\]

第二次透镜成实像，像距为

\[
v_3\approx35.6119\,\mathrm{cm}.
\]

另一方面，脸到透镜距离减去明视距离为

\[
u-25\approx35.6119\,\mathrm{cm},
\]

两者完全一致。逐级数值复核的意义在于：它不仅检查最终二次方程，还检查了反射后的传播方向和
\(u_3=d-v_2\) 的符号。

还可以用两个边界检验总映射。\(u\to0^+\) 时

\[
v_3\to420\,\mathrm{cm},
\]

而 \(u\to420^-\) 时 \(v_3\to0^+\)。这与实像区间 \((0,420)\)
的两个端点相互对应。

\newpage

\section{第 3 题}\label{ux7b2c-3-ux9898}

质量为 \(m\) 的小物体绕地球作半径为 \(1.5R\) 的圆周运动，其中 \(R\)
为地球半径，地球表面重力加速度为 \(g\)。忽略地球自转和空气阻力。

\begin{enumerate}
\def\labelenumi{\arabic{enumi}.}
\tightlist
\item
  求圆轨道周期；
\item
  从地球表面点 \(B\) 以第一宇宙速度、相对当地水平面仰角 \(\alpha\)
  发射质量为 \(nm\) 的第二个物体，使其椭圆轨道在点 \(C\) 与半径 \(1.5R\)
  的圆轨道相切。求 \(OB\) 与 \(x\) 轴的夹角 \(\theta_2\) 以及发射角
  \(\alpha\)；
\item
  为使两物体在 \(C\) 点相遇，应在质量为 \(m\) 的物体运行到点 \(A\)
  时发射第二物体。求 \(OA\) 与 \(x\) 轴的夹角 \(\theta_1\)；
\item
  两物体在 \(C\)
  点发生完全非弹性碰撞。要求碰后轨道与地球表面相切，求质量比 \(n\)。
\end{enumerate}

\begin{figure}[H]
\centering
\includegraphics[width=0.592\linewidth,height=\textheight,keepaspectratio,alt={图 3-1　圆轨道、转移椭圆及会合点的几何关系。}]{../figures/question-03.pdf}
\caption*{图 3-1　圆轨道、转移椭圆及会合点的几何关系。}
\end{figure}

\subsection{一、题型定位}\label{ux4e00ux9898ux578bux5b9aux4f4d-2}

椭圆轨道的几何中心与地心焦点不同，飞行时间所需面积应从焦点扫出。先求轨道参数和扫过面积，再处理与卫星迎面碰撞后的轨道变化。

\subsection{二、知识点}\label{ux4e8cux77e5ux8bc6ux70b9-2}

\begin{itemize}
\tightlist
\item
  \(GM=gR^2\)；
\item
  圆轨道速度与周期；
\item
  椭圆轨道的比机械能 \(\varepsilon=-GM/(2a)\)；
\item
  开普勒第二、第三定律；
\item
  切向对撞的一维动量守恒；
\item
  远地点、近地点与半长轴的关系。
\end{itemize}

\subsection{三、物理图像}\label{ux4e09ux7269ux7406ux56feux50cf-2}

圆轨道半径为 \(3R/2\)。第二物体从地表 \(r=R\)
处发射，以同样的第一宇宙速度大小进入一条半长轴为 \(R\) 的椭圆；点 \(C\)
是椭圆远地点，也是两条轨道的切点。第二物体从上方的 \(B\) 运动到 \(C\)
时切向速度向下，而第一物体在 \(C\) 的圆轨道速度向上，两者发生迎面碰撞。

\subsection{四、解题入口}\label{ux56dbux89e3ux9898ux5165ux53e3-2}

第二问先用能量求 \(C\)
点速度，再用角动量求发射角；椭圆几何由半长轴和焦距确定。第三问不能用``路程除速度''，而应使用焦点面积速度不变。第四问先由目标轨道的远、近地点求碰后速度，再列切向动量守恒。

\subsection{五、解析}\label{ux4e94ux89e3ux6790-2}

\begin{figure}[H]
\centering
\includegraphics[width=0.512\linewidth,height=\textheight,keepaspectratio,alt={图 3-2　椭圆扫面积计时与 C 点完全非弹性碰撞。}]{../figures/solution-03.pdf}
\caption*{图 3-2　椭圆扫面积计时与 \(C\) 点完全非弹性碰撞。}
\end{figure}

图 3-2 的阴影对应第 3 问：计时面积必须以地心焦点 \(O\)
为顶点，而不是以椭圆中心为顶点。图中 \(C\) 点的两支相反速度箭头对应第 4
问：在 \(C\) 点，圆轨道速度 \(v_c\) 与转移轨道速度 \(v_f\)
方向相反，碰后速度 \(u\) 与 \(v_c\) 同向，因此动量方程中两项异号。

\subsubsection{\texorpdfstring{1. 半径 \(1.5R\)
圆轨道的周期}{1. 半径 1.5R 圆轨道的周期}}\label{ux534aux5f84-1.5r-ux5706ux8f68ux9053ux7684ux5468ux671f}

地表重力给出

\[
GM=gR^2.
\]

设圆轨道半径为

\[
r_c=\frac{3R}{2}.
\]

万有引力提供向心力：

\[
\frac{GMm}{r_c^2}=m\frac{v_c^2}{r_c}.
\]

所以

\[
v_c=\sqrt{\frac{GM}{r_c}}
=\sqrt{\frac{2gR}{3}}.
\]

周期为

\[
\begin{aligned}
T_c
&=\frac{2\pi r_c}{v_c}\\
&=2\pi\sqrt{\frac{r_c^3}{GM}}\\
&=2\pi\sqrt{\frac{(3R/2)^3}{gR^2}}\\
&=3\pi\sqrt{\frac{3R}{2g}}.
\end{aligned}
\]

因此

\[
\boxed{T_c=3\pi\sqrt{\frac{3R}{2g}}}.
\]

\subsubsection{2.
转移椭圆及发射角}\label{ux8f6cux79fbux692dux5706ux53caux53d1ux5c04ux89d2}

地表第一宇宙速度为

\[
v_0=\sqrt{gR}=\sqrt{\frac{GM}{R}}.
\]

发射点的单位质量机械能为

\[
\varepsilon
=\frac{v_0^2}{2}-\frac{GM}{R}
=\frac{GM}{2R}-\frac{GM}{R}
=-\frac{GM}{2R}.
\]

椭圆轨道的比机械能为 \(-GM/(2a)\)，故

\[
a=R.
\]

点 \(C\) 是远地点，半径为 \(r_a=3R/2\)。椭圆近地点半径满足

\[
r_a+r_p=2a,
\]

因此

\[
r_p=2R-\frac{3R}{2}=\frac{R}{2}.
\]

焦距为

\[
c=\frac{r_a-r_p}{2}=\frac{R}{2},
\]

离心率与半短轴分别为

\[
e=\frac{c}{a}=\frac12,
\qquad
b=\sqrt{a^2-c^2}=\frac{\sqrt3}{2}R.
\]

在 \(C\) 点用机械能守恒：

\[
\frac{v_0^2}{2}-\frac{GM}{R}
=\frac{v_f^2}{2}-\frac{GM}{3R/2}.
\]

代入 \(v_0^2=GM/R\)：

\[
-\frac{GM}{2R}
=\frac{v_f^2}{2}-\frac{2GM}{3R}.
\]

于是

\[
\frac{v_f^2}{2}=\frac{GM}{6R},
\qquad
\boxed{v_f=\sqrt{\frac{gR}{3}}}.
\]

发射时只有水平分速度对地心产生角动量。以单位质量写角动量守恒：

\[
Rv_0\cos\alpha=\frac{3R}{2}v_f.
\]

所以

\[
\cos\alpha
=\frac{3v_f}{2v_0}
=\frac{3}{2\sqrt3}
=\frac{\sqrt3}{2}.
\]

因此

\[
\boxed{\alpha=30^\circ}.
\]

再求发射点的极角。以远地点方向为极轴，椭圆焦点极坐标式可写为

\[
r=\frac{p}{1-e\cos\theta},
\]

其中

\[
p=\frac{b^2}{a}=\frac{3R}{4}.
\]

在地表交点 \(B\)，\(r=R\)，故

\[
R=\frac{3R/4}{1-(1/2)\cos\theta_2}.
\]

约去 \(R\)：

\[
1-\frac12\cos\theta_2=\frac34.
\]

于是

\[
\cos\theta_2=\frac12,
\qquad
\boxed{\theta_2=60^\circ}.
\]

\subsubsection{3.
会合所需的初始圆轨道角位置}\label{ux4f1aux5408ux6240ux9700ux7684ux521dux59cbux5706ux8f68ux9053ux89d2ux4f4dux7f6e}

第二物体从 \(B\) 到 \(C\) 的时间由开普勒第二定律求得。椭圆中心记为
\(O'\)，则 \(O'O=R/2\)。点 \(B\)
恰好位于椭圆短轴端点，因为相对椭圆中心的横坐标为零、纵坐标为 \(b\)。

椭圆中心扇形 \(BO'C\) 是四分之一椭圆，其面积为

\[
S_{BO'C}=\frac{\pi ab}{4}.
\]

三角形 \(BO'O\) 的面积为

\[
S_{BO'O}=\frac12\cdot\frac{R}{2}\cdot b
=\frac{ab}{4},
\]

因为 \(a=R\)。因此焦点 \(O\) 的扫过面积为

\[
S_{BOC}=\frac{ab}{4}(\pi+1).
\]

椭圆总面积是 \(\pi ab\)。设半长轴为 \(R\) 的椭圆周期为 \(T_e\)，则

\[
t_{BC}=\frac{S_{BOC}}{\pi ab}T_e
=\frac{\pi+1}{4\pi}T_e.
\]

质量为 \(m\) 的物体在半径 \(3R/2\) 的圆轨道上，从 \(A\) 到 \(C\)
转过角度 \(\theta_1\)，所需时间为

\[
t_{AC}=\frac{\theta_1}{2\pi}T_c.
\]

会合要求 \(t_{AC}=t_{BC}\)。开普勒第三定律给出

\[
\left(\frac{T_e}{T_c}\right)^2
=\left(\frac{R}{3R/2}\right)^3
=\left(\frac23\right)^3.
\]

所以

\[
\frac{T_e}{T_c}=\left(\frac23\right)^{3/2}.
\]

代入时间相等条件：

\[
\frac{\theta_1}{2\pi}T_c
=\frac{\pi+1}{4\pi}T_e.
\]

约去周期后

\[
\begin{aligned}
\theta_1
&=\frac{\pi+1}{2}\left(\frac23\right)^{3/2}\\
&=\frac{\sqrt6}{9}(\pi+1).
\end{aligned}
\]

因此

\[
\boxed{\theta_1=\frac{\sqrt6}{9}(\pi+1)\ \mathrm{rad}}.
\]

\subsubsection{4.
完全非弹性碰撞后的质量比}\label{ux5b8cux5168ux975eux5f39ux6027ux78b0ux649eux540eux7684ux8d28ux91cfux6bd4}

碰后轨道在 \(C\) 点的远地点半径为
\(3R/2\)，并与地球表面相切，所以近地点半径为 \(R\)。新椭圆的半长轴为

\[
a'=\frac{3R/2+R}{2}=\frac{5R}{4}.
\]

设碰后在 \(C\) 点的速度为 \(u\)。由比机械能：

\[
\frac{u^2}{2}-\frac{GM}{3R/2}
=-\frac{GM}{2a'}.
\]

代入 \(a'=5R/4\)：

\[
\frac{u^2}{2}-\frac{2GM}{3R}
=-\frac{2GM}{5R}.
\]

所以

\[
\frac{u^2}{2}
=\left(\frac23-\frac25\right)\frac{GM}{R}
=\frac{4GM}{15R}.
\]

由此

\[
\boxed{u=\sqrt{\frac{8gR}{15}}}.
\]

在 \(C\) 点，圆轨道物体速度为

\[
v_c=\sqrt{\frac{2gR}{3}},
\]

转移椭圆物体速度为

\[
v_f=\sqrt{\frac{gR}{3}}.
\]

取圆轨道物体的速度方向为正。两物体迎面运动，完全非弹性碰撞后粘在一起，切向动量守恒：

\[
mv_c-nm v_f=(n+1)m u.
\]

约去 \(m\) 并整理：

\[
n(v_f+u)=v_c-u,
\]

\[
n=\frac{v_c-u}{v_f+u}.
\]

把各速度统一写成 \(\sqrt{gR/15}\) 的倍数：

\[
v_c=\sqrt{10}\sqrt{\frac{gR}{15}},
\quad
v_f=\sqrt5\sqrt{\frac{gR}{15}},
\quad
u=2\sqrt2\sqrt{\frac{gR}{15}}.
\]

因此

\[
\boxed{
n=\frac{\sqrt{10}-2\sqrt2}{\sqrt5+2\sqrt2}
\approx0.0660
}.
\]

\subsubsection{5.
三个速度尺度的相互检验}\label{ux4e09ux4e2aux901fux5ea6ux5c3aux5ea6ux7684ux76f8ux4e92ux68c0ux9a8c}

在 \(C\) 点，三种速度的平方依次为

\[
v_c^2=\frac{10}{15}gR,
\qquad
v_f^2=\frac{5}{15}gR,
\qquad
u^2=\frac{8}{15}gR.
\]

所以

\[
v_c>u>v_f.
\]

圆轨道物体的正向动量必须先抵消第二物体的反向动量，再给粘连体留下正向速度
\(u\)。因此 \(n\) 应当是一个小于 \(1\) 的正数，计算值 \(0.0660\)
与这个定性判断一致。

碰撞前后的机械能不能守恒。以 \(\sqrt{gR/15}\) 为速度单位，碰前动能为

\[
K_i=\frac12m\frac{gR}{15}(10+5n),
\]

碰后动能为

\[
K_f=\frac12(1+n)m\frac{gR}{15}\cdot8.
\]

用动量方程得到的 \(n\) 代入后有
\(K_f<K_i\)，差值转化为内能，符合完全非弹性碰撞。

会合时间也可用偏近点角复核。点 \(B\) 是椭圆短轴端点，对应偏近点角
\(E_B=\pi/2\)，远地点 \(C\) 对应
\(E_C=\pi\)。开普勒方程中的平均近点角增量为

\[
\Delta M=(E_C-e\sin E_C)-(E_B-e\sin E_B)
=\frac{\pi+1}{2}.
\]

除以 \(2\pi\) 后同样得到

\[
\frac{t_{BC}}{T_e}=\frac{\pi+1}{4\pi},
\]

与面积法完全一致。

\newpage

\section{第 4 题}\label{ux7b2c-4-ux9898}

一根硬质电阻丝制成``8''字形闭合回路，左右两部分均为圆，直径分别为
\(D\)、\(d\)，且 \(D>d\)。两圆在接触点处绝缘，回路总电阻为
\(R\)。回路位于光滑水平面内，左侧无磁场，右侧存在竖直向下、磁感应强度为
\(B\) 的匀强磁场。回路以垂直磁场边界的初速度 \(v_0\)
向右运动。记回路总质量为 \(m\)。

\begin{enumerate}
\def\labelenumi{\arabic{enumi}.}
\tightlist
\item
  若回路最终完全进入磁场，求此过程中流过回路任一点的电荷量绝对值；
\item
  以回路最右端进入磁场的距离为 \(x\)，求运动速度 \(v(x)\)；
\item
  求使回路能够全部进入磁场区域所需的最小初速度。
\end{enumerate}

\begin{figure}[H]
\centering
\includegraphics[width=0.808\linewidth,height=\textheight,keepaspectratio,alt={图 4-1　``8''字形电阻丝进入有界匀强磁场。}]{../figures/question-04.pdf}
\caption*{图 4-1　``8''字形电阻丝进入有界匀强磁场。}
\end{figure}

\subsection{一、题型定位}\label{ux4e00ux9898ux578bux5b9aux4f4d-3}

两个相接圆环先后进入磁场，感应电流在两个阶段反向，磁阻力却都阻碍运动。把净感应电荷、累计通过电荷和耗散能量分别计算。

\subsection{二、知识点}\label{ux4e8cux77e5ux8bc6ux70b9-3}

\begin{itemize}
\tightlist
\item
  净感应电荷 \(dQ=-d\Phi/R\) 与累计通电量 \(dQ_{\rm abs}=|d\Phi|/R\)；
\item
  ``8''字回路两叶的绕向相反；
\item
  动生电动势 \(\mathcal E=B\ell v\)；
\item
  安培力 \(F=B I\ell\)；
\item
  用 \(dv/dt=v\,dv/dx\) 消去时间；
\item
  圆被直线截得的弦长。
\end{itemize}

\subsection{三、物理图像}\label{ux4e09ux7269ux7406ux56feux50cf-3}

取 \(x=0\) 为小圆最右端刚接触磁场边界。\(0<x<d\)
时只有小圆被边界切割；\(d<x<d+D\)
时小圆已完全进入，边界开始切割大圆；\(x>d+D\)
时整个回路进入匀强磁场，磁通量不再变化，感应电流消失。

题面若不提供质量，就无法由安培力求速度变化。这里明确记总质量为
\(m\)，与手写解答中使用的动力学参数一致。

\subsection{四、解题入口}\label{ux56dbux89e3ux9898ux5165ux53e3-3}

第一问直接对法拉第定律积分，不必求瞬时电流。第二问要先写弦长
\(\ell(x)\)，再把安培力写成 \(B^2\ell^2v/R\)。利用 \(dx/dt=v\)
后，速度方程中的 \(v\) 会约掉，速度损失只剩对 \(\ell^2\) 的空间积分。

\subsection{五、解析}\label{ux4e94ux89e3ux6790-3}

\begin{figure}[H]
\centering
\includegraphics[width=0.9\linewidth,height=\textheight,keepaspectratio,alt={图 4-2　小圆、大圆依次被磁场边界切割的三个阶段。}]{../figures/solution-04.pdf}
\caption*{图 4-2　小圆、大圆依次被磁场边界切割的三个阶段。}
\end{figure}

图 4-2 左图把小圆最右端进入边界的距离记为 \(x\)，边界截得的弦长记为
\(\ell_d\)；中图把大圆最右端进入边界的距离记为 \(x-d\)，对应弦长记为
\(\ell_D\)。正文再把它们写成函数
\(\ell_d(x)\)、\(\ell_D(x-d)\)。这两个弦长同时决定动生电动势 \(B\ell v\)
和安培力 \(BI\ell\)，所以速度方程中出现的是 \(\ell^2\)。

\subsubsection{1.
净感应电荷与累计通电量}\label{ux51c0ux611fux5e94ux7535ux8377ux4e0eux7d2fux8ba1ux901aux7535ux91cf}

固定回路绕行正方向，\(I=-\dot\Phi/R\)。整个进入过程的净电荷为

\[Q_{\rm net}=\int I\,dt=-\frac{\Phi_f-\Phi_i}{R},\qquad \boxed{|Q_{\rm net}|=\frac{B\pi(D^2-d^2)}{4R}}.\]

这对应原参考解的``电荷量绝对值''。但两叶绕向相反，小圆与大圆进入时电流反向，因此累计流过的电荷大小是另一量：

\[Q_{\rm abs}=\int|I|\,dt=\frac1R\int|d\Phi|=\boxed{\frac{B\pi(D^2+d^2)}{4R}}.\]

不能从 \(dq=|d\Phi|/R\)
积分得到面积相减。题面``电荷量绝对值''应区分是先积分后取绝对值，还是累计两个方向的通电量；两种解释及结果在这里一并列出。后续速度损失取决于弦长平方，两阶段总是相加。

\subsubsection{\texorpdfstring{2. 小圆进入阶段
\(0\le x\le d\)}{2. 小圆进入阶段 0\textbackslash le x\textbackslash le d}}\label{ux5c0fux5706ux8fdbux5165ux9636ux6bb5-0le-xle-d}

小圆半径为 \(d/2\)。边界到小圆最右端的距离为
\(x\)，因此边界相对圆心的横坐标为 \(d/2-x\)。由圆方程，切割弦长为

\[
\begin{aligned}
\ell_d(x)
&=2\sqrt{\left(\frac d2\right)^2-\left(\frac d2-x\right)^2}\\
&=2\sqrt{dx-x^2}.
\end{aligned}
\]

动生电动势的绝对值为

\[
|\mathcal E|=B\ell_d v,
\]

电流为

\[
I=\frac{B\ell_d v}{R}.
\]

磁场对处在边界内的有效直导线段产生与运动相反的合力，大小为

\[
F_m=B I\ell_d
=\frac{B^2\ell_d^2}{R}v.
\]

牛顿第二定律为

\[
m\frac{dv}{dt}=-\frac{B^2\ell_d^2}{R}v.
\]

利用

\[
\frac{dv}{dt}=\frac{dv}{dx}\frac{dx}{dt}=v\frac{dv}{dx},
\]

且运动过程中 \(v>0\)，约去 \(v\)：

\[
m\frac{dv}{dx}=-\frac{B^2\ell_d^2}{R}.
\]

代入 \(\ell_d^2=4(dx-x^2)\)：

\[
\frac{dv}{dx}=-\frac{4B^2}{mR}(dx-x^2).
\]

从 \(0\) 积分到 \(x\)：

\[
\begin{aligned}
v(x)-v_0
&=-\frac{4B^2}{mR}\int_0^x(d\xi-\xi^2)\,d\xi\\
&=-\frac{4B^2}{mR}
\left(\frac{d x^2}{2}-\frac{x^3}{3}\right).
\end{aligned}
\]

所以

\[
\boxed{
v(x)=v_0-\frac{4B^2}{mR}
\left(\frac{d x^2}{2}-\frac{x^3}{3}\right),
\quad 0\le x\le d
}.
\]

当 \(x=d\) 时，速度为

\[
v(d)=v_0-\frac{2B^2d^3}{3mR}.
\]

\subsubsection{\texorpdfstring{3. 大圆进入阶段
\(d\le x\le d+D\)}{3. 大圆进入阶段 d\textbackslash le x\textbackslash le d+D}}\label{ux5927ux5706ux8fdbux5165ux9636ux6bb5-dle-xle-dd}

令

\[
X=x-d.
\]

此时 \(X\) 是大圆最右端进入磁场的距离，范围为
\(0\le X\le D\)。大圆切割弦长为

\[
\ell_D(X)=2\sqrt{DX-X^2}.
\]

与上一阶段完全同理：

\[
\frac{dv}{dX}=-\frac{4B^2}{mR}(DX-X^2).
\]

从 \(X=0\) 积分到 \(X=x-d\)，并以 \(v(d)\) 为初值：

\[
\begin{aligned}
v(x)
&=v(d)-\frac{4B^2}{mR}
\left[\frac{D(x-d)^2}{2}-\frac{(x-d)^3}{3}\right]\\
&=v_0-\frac{2B^2d^3}{3mR}
-\frac{4B^2}{mR}
\left[\frac{D(x-d)^2}{2}-\frac{(x-d)^3}{3}\right].
\end{aligned}
\]

即

\[
\boxed{
v(x)=v_0-\frac{2B^2d^3}{3mR}
-\frac{4B^2}{mR}
\left[\frac{D(x-d)^2}{2}-\frac{(x-d)^3}{3}\right],
\quad d\le x\le d+D
}.
\]

\subsubsection{4.
完全进入磁场以后}\label{ux5b8cux5168ux8fdbux5165ux78c1ux573aux4ee5ux540e}

在 \(x=d+D\) 处，第二阶段的附加速度损失为

\[
\frac{4B^2}{mR}\left(\frac{D^3}{2}-\frac{D^3}{3}\right)
=\frac{2B^2D^3}{3mR}.
\]

所以完全进入后速度保持为

\[
\boxed{
v_\infty=v_0-\frac{2B^2(D^3+d^3)}{3mR}
}.
\]

把三段合并，可写成

\[
v(x)=
\begin{cases}
v_0-\dfrac{4B^2}{mR}
\left(\dfrac{dx^2}{2}-\dfrac{x^3}{3}\right),
&0\le x\le d,\\[8pt]
v_0-\dfrac{2B^2d^3}{3mR}
-\dfrac{4B^2}{mR}
\left[\dfrac{D(x-d)^2}{2}-\dfrac{(x-d)^3}{3}\right],
&d\le x\le d+D,\\[8pt]
v_0-\dfrac{2B^2(D^3+d^3)}{3mR},
&x\ge d+D.
\end{cases}
\]

\subsubsection{5. 最小初速度}\label{ux6700ux5c0fux521dux901fux5ea6}

要使回路不在进入过程中停下，必须有

\[
v_\infty\ge0.
\]

因此初速度的下界为

\[
\boxed{
v_{0,\min}=\frac{2B^2(D^3+d^3)}{3mR}
}.
\]

若要求在有限时间内完全越过边界，应取严格大于此值；等号给出能够完成进入的临界下界。

\subsubsection{6.
用能量耗散交叉验证速度损失}\label{ux7528ux80fdux91cfux8017ux6563ux4ea4ux53c9ux9a8cux8bc1ux901fux5ea6ux635fux5931}

在某个进入阶段，焦耳热功率为

\[
P_J=I^2R
=\frac{B^2\ell^2v^2}{R}.
\]

机械能减少率为

\[
-\frac{d}{dt}\left(\frac12mv^2\right)
=-mv\frac{dv}{dt}.
\]

能量守恒要求机械能的减少率等于焦耳热功率：

\[
-mv\frac{dv}{dt}
=\frac{B^2\ell^2v^2}{R}.
\]

在 \(v>0\) 时约去一个 \(v\)，再用 \(dx/dt=v\)：

\[
m\frac{dv}{dx}=-\frac{B^2\ell^2}{R}.
\]

这正是由安培力得到的空间微分方程，说明磁阻力法与能量法相互一致。

还可作量纲检查：

\[
\frac{B^2d^3}{mR}
\]

的量纲为速度。因为 \(B\ell v\) 是电压，\(B^2\ell^2v/R\)
是力，再乘一个长度并除以质量与速度，最终得到速度损失。若误把结果写成与
\(d^2\) 成正比，量纲就不能成立。

临界初速度取等号时，\(v\)
在最后边界位置降为零。由于到达终点前速度越来越小，严格完成进入需要无限逼近；实际要求有限时间完全进入时，应保留一个正的速度余量。

\textbf{最小速度的等号对应无限时间}

在最后进入的阶段，令剩余距离 \(\delta=d+D-x\)。边界弦长平方约为
\(4D\delta\)，若初速度恰等于正文阈值，剩余速度为
\(v\simeq2B^2D\delta^2/(mR)\)。于是 \(dt\propto d\delta/\delta^2\)，到达
\(\delta=0\)
所需时间发散。阈值是可完整进入的速度下确界；要求有限时间完全越过边界，应严格大于它。

\newpage

\section{第 5 题}\label{ux7b2c-5-ux9898}

三个固定点 \(A\)、\(C\)、\(E\) 满足：\(AC\) 为水平线段，长度为
\(d\)；\(EC\) 为竖直线段，长度为 \(h\)。可伸缩杆 \(AB\) 绕 \(A\)
转动。刚性折杆 \(BCD\) 绕 \(C\) 转动，其中 \(BC=a\)、\(CD=h\)，且
\(BC\perp CD\)。固定点 \(E\) 处有一可转动滑套，长度为 \(2h\) 的杆 \(DF\)
穿过滑套并可沿套滑动，杆端 \(D\) 与折杆铰接。初始时 \(DF\) 竖直；随着
\(AB\) 伸长，\(DF\) 转动，端点 \(F\) 始终位于桌面上方。

\begin{enumerate}
\def\labelenumi{\arabic{enumi}.}
\tightlist
\item
  以 \(E\) 为原点，求端点 \(F\) 的轨迹；
\item
  当 \(DF\) 与桌面夹角为 \(\theta\)、\(AB\) 的伸长速度为 \(v\) 时，求
  \(F\) 点的速度。
\end{enumerate}

\begin{figure}[H]
\centering
\includegraphics[width=0.442\linewidth,height=\textheight,keepaspectratio,alt={图 5-1　可伸缩杆、刚性折杆与滑动套组成的平面机构。}]{../figures/question-05.pdf}
\caption*{图 5-1　可伸缩杆、刚性折杆与滑动套组成的平面机构。}
\end{figure}

\subsection{一、题型定位}\label{ux4e00ux9898ux578bux5b9aux4f4d-4}

杆和滑环的约束给出一个选定分支的端点轨迹。先用驱动杆长求角速度，再按图中的坐标分解端点速度，检查机构接近死点时的分母。

\subsection{二、知识点}\label{ux4e8cux77e5ux8bc6ux70b9-4}

\begin{itemize}
\tightlist
\item
  等腰三角形中的弦长关系；
\item
  极坐标形式的轨迹方程；
\item
  余弦定理与约束微分；
\item
  复合函数求导；
\item
  平面速度的直角坐标分量和极坐标分量。
\end{itemize}

\subsection{三、物理图像}\label{ux4e09ux7269ux7406ux56feux50cf-4}

在竖直平面内取 \(E\) 为原点，沿桌面向左为 \(x\) 轴正方向，竖直向上为
\(y\) 轴。杆 \(DF\) 通过固定滑套 \(E\)，故 \(D,E,F\) 始终共线。由于
\(EC=CD=h\)，三角形 \(ECD\) 是等腰三角形。给定杆的倾角 \(\theta\)
后，\(ED\) 以及剩余长度 \(EF=DF-ED\) 都被唯一确定。

\subsection{四、解题入口}\label{ux56dbux89e3ux9898ux5165ux53e3-4}

第一问先在三角形 \(ECD\) 中求 \(ED\)，再把 \(F\) 写成极坐标
\(r=EF\)、极角 \(\theta\)。第二问用余弦定理把 \(AB\) 的长度写成
\(\theta\) 的函数，微分得到 \(\dot\theta\)；最后对 \(F\) 的坐标求导。

\subsection{五、解析}\label{ux4e94ux89e3ux6790-4}

\begin{figure}[H]
\centering
\includegraphics[width=0.438\linewidth,height=\textheight,keepaspectratio,alt={图 5-2　以 E 为原点建立坐标系并分解端点速度。}]{../figures/solution-05.pdf}
\caption*{图 5-2　以 \(E\) 为原点建立坐标系并分解端点速度。}
\end{figure}

图 5-2 同时服务于轨迹和速度两部分。点 \(F\) 的极径为 \(r=EF\)，极角为
\(\theta\)；速度不是沿杆方向的单一分量，而是径向分量 \(\dot r\,\bm e_r\)
与横向分量 \(r\dot\theta\,\bm e_\theta\)
的矢量和。后面的直角坐标求导与极坐标分解应得到同一个 \(\bm v_F\)。

\subsubsection{\texorpdfstring{1. 求 \(ED\) 与
\(EF\)}{1. 求 ED 与 EF}}\label{ux6c42-ed-ux4e0e-ef}

设 \(DF\) 与水平桌面的夹角为 \(\theta\)。在三角形 \(ECD\) 中，

\[
EC=CD=h.
\]

从 \(E\) 向 \(C\) 的方向竖直向下，\(ED\) 与水平线夹角为 \(\theta\)，所以
\(\angle CED=90^\circ-\theta\)。等腰三角形中另一个底角相等，故

\[
\angle EDC=90^\circ-\theta.
\]

顶角为

\[
\angle ECD=180^\circ-2(90^\circ-\theta)=2\theta.
\]

等腰三角形底边为

\[
ED=2h\sin\theta.
\]

由于整根杆长为 \(DF=2h\)，且滑套 \(E\) 位于 \(D\)、\(F\) 之间，因此

\[
EF=DF-ED=2h-2h\sin\theta.
\]

记

\[
r=EF,
\]

便得到

\[
\boxed{r=2h(1-\sin\theta)}.
\]

\subsubsection{\texorpdfstring{2. 端点 \(F\)
的轨迹}{2. 端点 F 的轨迹}}\label{ux7aefux70b9-f-ux7684ux8f68ux8ff9}

以从 \(E\) 指向 \(F\) 的方向作为极径方向，则

\[
x=r\cos\theta,
\qquad
y=r\sin\theta.
\]

代入 \(r=2h(1-\sin\theta)\)：

\[
x=2h(1-\sin\theta)\cos\theta,
\]

\[
y=2h(1-\sin\theta)\sin\theta.
\]

要消去 \(\theta\)，先利用

\[
r=\sqrt{x^2+y^2},
\qquad
\sin\theta=\frac{y}{r}.
\]

代入极径方程：

\[
r=2h\left(1-\frac{y}{r}\right).
\]

两边乘以 \(r\)：

\[
r^2=2hr-2hy.
\]

再用 \(r^2=x^2+y^2\)：

\[
x^2+y^2=2h\sqrt{x^2+y^2}-2hy.
\]

因此轨迹的隐式方程为

\[
\boxed{
x^2+y^2+2hy-2h\sqrt{x^2+y^2}=0
}.
\]

轨迹方程只描述机构实际可达的那一段，还应同时满足杆、滑套和铰链给出的角度范围，不能把代数曲线的全部分支都当作真实轨迹。

\subsubsection{\texorpdfstring{3. 用 \(AB\) 的伸长速度求
\(\dot\theta\)}{3. 用 AB 的伸长速度求 \textbackslash dot\textbackslash theta}}\label{ux7528-ab-ux7684ux4f38ux957fux901fux5ea6ux6c42-dottheta}

记

\[
\ell=AB.
\]

由于刚性折杆满足 \(BC\perp CD\)，而 \(CD\) 的方向由 \(\theta\)
决定，可得 \(BC\) 与固定水平线 \(AC\) 的夹角

\[
\varphi=\pi-2\theta.
\]

在三角形 \(ABC\) 中，\(AC=d\)、\(BC=a\)。由余弦定理：

\[
\ell^2=a^2+d^2-2ad\cos\varphi.
\]

利用 \(\cos(\pi-2\theta)=-\cos2\theta\)：

\[
\boxed{
\ell^2=a^2+d^2+2ad\cos2\theta
}.
\]

已知杆的伸长速度为

\[
\dot\ell=v.
\]

对约束式两边关于时间求导：

\[
2\ell\dot\ell
=2ad\frac{d}{dt}(\cos2\theta).
\]

而

\[
\frac{d}{dt}(\cos2\theta)
=-2\sin2\theta\,\dot\theta.
\]

所以

\[
2\ell v=-4ad\sin2\theta\,\dot\theta.
\]

解得

\[
\boxed{
\dot\theta
=-\frac{\sqrt{a^2+d^2+2ad\cos2\theta}}
{2ad\sin2\theta}\,v
}.
\]

负号表明在图示支路中，\(AB\) 伸长时 \(\theta\) 减小。

\subsubsection{\texorpdfstring{4. 求 \(F\)
点速度分量}{4. 求 F 点速度分量}}\label{ux6c42-f-ux70b9ux901fux5ea6ux5206ux91cf}

位置坐标为

\[
x_F=2h(1-\sin\theta)\cos\theta,
\]

\[
y_F=2h(1-\sin\theta)\sin\theta.
\]

先对 \(x_F\) 求导：

\[
\begin{aligned}
\dot x_F
&=2h\frac{d}{dt}\left[(1-\sin\theta)\cos\theta\right]\\
&=2h\left[-\cos^2\theta-(1-\sin\theta)\sin\theta\right]\dot\theta\\
&=-2h\left(\cos2\theta+\sin\theta\right)\dot\theta.
\end{aligned}
\]

代入 \(\dot\theta\)：

\[
\boxed{
v_{Fx}
=\frac{h\sqrt{a^2+d^2+2ad\cos2\theta}}
{ad\sin2\theta}
(\sin\theta+\cos2\theta)v
}.
\]

再对 \(y_F\) 求导：

\[
\begin{aligned}
\dot y_F
&=2h\frac{d}{dt}\left[(1-\sin\theta)\sin\theta\right]\\
&=2h\left[-\cos\theta\sin\theta+(1-\sin\theta)\cos\theta\right]\dot\theta\\
&=2h\cos\theta(1-2\sin\theta)\dot\theta.
\end{aligned}
\]

代入 \(\dot\theta\) 后：

\[
\boxed{
v_{Fy}
=\frac{h\sqrt{a^2+d^2+2ad\cos2\theta}}
{ad\sin2\theta}
(2\sin\theta-1)\cos\theta\,v
}.
\]

\subsubsection{5.
速度大小的简化}\label{ux901fux5ea6ux5927ux5c0fux7684ux7b80ux5316}

也可以直接使用极坐标速度。由

\[
r=2h(1-\sin\theta)
\]

得

\[
\dot r=-2h\cos\theta\,\dot\theta.
\]

速度平方为

\[
v_F^2=\dot r^2+r^2\dot\theta^2.
\]

代入 \(r\)、\(\dot r\)：

\[
\begin{aligned}
v_F^2
&=4h^2\left[\cos^2\theta+(1-\sin\theta)^2\right]\dot\theta^2\\
&=4h^2(2-2\sin\theta)\dot\theta^2.
\end{aligned}
\]

因此

\[
v_F=2h\sqrt{2-2\sin\theta}\,|\dot\theta|.
\]

在图示角度范围内代入 \(|\dot\theta|\)：

\[
\boxed{
v_F=
\frac{h\sqrt{a^2+d^2+2ad\cos2\theta}}
{ad\sin2\theta}
\sqrt{2-2\sin\theta}\,v
}.
\]

这个结果和由 \(v_{Fx}^2+v_{Fy}^2\) 计算所得完全一致。

\subsubsection{6.
速度分量的代数复核与奇异位置}\label{ux901fux5ea6ux5206ux91cfux7684ux4ee3ux6570ux590dux6838ux4e0eux5947ux5f02ux4f4dux7f6e}

由极坐标单位矢量

\[
\bm e_r=(\cos\theta,\sin\theta),
\qquad
\bm e_\theta=(-\sin\theta,\cos\theta),
\]

端点速度还可写成

\[
\bm v_F=\dot r\,\bm e_r+r\dot\theta\,\bm e_\theta.
\]

代入 \(\dot r=-2h\cos\theta\dot\theta\)、\(r=2h(1-\sin\theta)\)，其
\(x\) 分量为

\[
\begin{aligned}
v_{Fx}
&=-2h\cos^2\theta\dot\theta
-2h(1-\sin\theta)\sin\theta\dot\theta\\
&=-2h(\cos2\theta+\sin\theta)\dot\theta,
\end{aligned}
\]

而 \(y\) 分量为

\[
\begin{aligned}
v_{Fy}
&=-2h\sin\theta\cos\theta\dot\theta
+2h(1-\sin\theta)\cos\theta\dot\theta\\
&=2h(1-2\sin\theta)\cos\theta\dot\theta.
\end{aligned}
\]

这与直接对直角坐标求导完全相同。

结果中的 \(\sin2\theta\) 出现在分母，表示机构在 \(\sin2\theta=0\)
附近存在运动学奇异性：驱动杆长度对 \(\theta\)
的一阶变化趋于零，为维持给定的有限伸长速度 \(v\)，理论角速度
\(|\dot\theta|\)
会变得很大。真实机构会受到间隙、杆件惯性和有限驱动力限制，因此该理想表达式只适用于远离奇异位形的区间。

\textbf{坐标分支与速度箭头}

原图 \(F\) 在 \(E\) 左上方，故这里选 \(+x\) 向左、\(+y\)
向上，\(\theta\) 从向左的水平轴量起，才能保持
\(x=r\cos\theta,y=r\sin\theta\)。若坚持 \(+x\) 向右，则所有 \(x\) 坐标与
\(v_x\) 都应反号，速率大小不变。

在图示伸长阶段，\(\dot\theta<0\)。因此 \(r\dot\theta\bm e_\theta\)
指向左下，\(\dot r\bm e_r\)
指向左上，合速度应按数值矢量相加；不能仅看两条未带符号的基矢决定箭头。折杆
\(BC\) 与 \(CD\) 始终垂直，解析图按同一 \(\theta\)
建出两段，避免示意图把余弦定理所需角度画错。

\newpage

\section{第 6 题}\label{ux7b2c-6-ux9898}

半径为 \(R\) 的半球形带电体球心固定在原点，对称轴与 \(x\)
轴重合，半球位于 \(x\le0\) 一侧。

\begin{enumerate}
\def\labelenumi{\arabic{enumi}.}
\tightlist
\item
  若仅圆形边缘均匀带电，线电荷密度为 \(\lambda\)，求 \(+x\)
  轴上电场强度的最大值及取得最大值的位置；
\item
  若仅半球弧面均匀带电，面电荷密度为 \(\sigma\)，求原点处的电场强度；
\item
  接第2问，求 \(+x\) 轴上任意点 \(x_0\) 处的电场强度。
\end{enumerate}

\begin{figure}[H]
\centering
\includegraphics[width=0.793\linewidth,height=\textheight,keepaspectratio,alt={图 6-1　半球弧面、圆形边缘及对称轴的空间关系。}]{../figures/question-06.pdf}
\caption*{图 6-1　半球弧面、圆形边缘及对称轴的空间关系。}
\end{figure}

\subsection{一、题型定位}\label{ux4e00ux9898ux578bux5b9aux4f4d-5}

圆环和半球都具有轴对称性，横向电场分量成对抵消。先积轴向分量，再检查圆环场的峰值、半球原点场及远场极限。

\subsection{二、知识点}\label{ux4e8cux77e5ux8bc6ux70b9-5}

\begin{itemize}
\tightlist
\item
  圆环轴线电场；
\item
  轴对称面元 \(dS=2\pi R^2\sin\theta\,d\theta\)；
\item
  只保留横向抵消后的轴向分量；
\item
  换元积分；
\item
  近场、原点和远场极限检验。
\end{itemize}

\subsection{三、物理图像}\label{ux4e09ux7269ux7406ux56feux50cf-5}

圆形边缘位于平面 \(x=0\)，半径为 \(R\)。半球弧面向 \(x<0\) 一侧鼓出。对
\(x\) 轴上的场点，任一绕轴圆周上的横向电场分量成对抵消，因此总场一定沿
\(x\) 轴。

\subsection{四、解题入口}\label{ux56dbux89e3ux9898ux5165ux53e3-5}

第一问可把边缘看成总电荷 \(Q=2\pi R\lambda\)
的圆环。第二、三问把半球弧面分割成与 \(x\)
轴同轴的窄圆环；每个窄环的场只需取轴向分量，原来的二重面积分便化成关于极角的一重积分。

\subsection{五、解析}\label{ux4e94ux89e3ux6790-5}

\begin{figure}[H]
\centering
\includegraphics[width=0.9\linewidth,height=\textheight,keepaspectratio,alt={图 6-2　轴对称电场的几何分解：（a）圆环电荷元；（b）半球面元的子午截面。}]{../figures/solution-06.pdf}
\caption*{图
6-2　轴对称电场的几何分解：（a）圆环电荷元；（b）半球面元的子午截面。}
\end{figure}

图 6-2 给出下面三问共同使用的几何关系。左图用于第 1
问：由直角三角形直接读出电荷元到场点的距离 \(s\) 以及轴向投影
\(dE_x\)。右图用于第 2、3 问：图中短弧绕 \(x\) 轴一周形成窄纬圈面元
\(dS\)；先写出 \(dS\)、\(s\)
和电场的轴向分量，再对极角积分。图中半球明确位于 \(x\le 0\)
一侧，因此面元的轴坐标为负值，这是第 3 问投影式中出现加号的原因。

记库仑常量为

\[
k_e=\frac{1}{4\pi\varepsilon_0}.
\]

\subsubsection{1.
圆形边缘带电时的轴线场}\label{ux5706ux5f62ux8fb9ux7f18ux5e26ux7535ux65f6ux7684ux8f74ux7ebfux573a}

圆环总电荷为

\[
Q=2\pi R\lambda.
\]

先不直接套用圆环轴线场公式。由图 6-2(a)，圆环上任一电荷元 \(dq\) 到场点
\(P\) 的距离均为

\[
s=\sqrt{R^2+x^2}.
\]

该电荷元产生的电场大小为

\[
dE=k_e\frac{dq}{s^2}.
\]

电场方向沿 \(dqP\)。图中 \(\alpha\) 为 \(d\bm E\) 与 \(+x\)
轴的夹角，因此

\[
\cos\alpha=\frac{x}{s}.
\]

圆环上关于轴线对称的电荷元，其垂直于 \(x\)
轴的电场分量成对抵消；轴向分量方向相同。于是一个电荷元贡献的轴向分量为

\[
\begin{aligned}
dE_x
&=dE\cos\alpha\\
&=k_e\frac{dq}{s^2}\frac{x}{s}\\
&=k_e\frac{x\,dq}{(R^2+x^2)^{3/2}}.
\end{aligned}
\]

式中 \(x\)、\(R\) 和 \(s\) 对整条圆环都相同，只有 \(dq\)
随积分位置变化。因此

\[
\begin{aligned}
E_x(x)
&=\oint dE_x\\
&=k_e\frac{x}{(R^2+x^2)^{3/2}}\oint dq\\
&=k_e\frac{Qx}{(R^2+x^2)^{3/2}}.
\end{aligned}
\]

代入 \(Q\)：

\[
E(x)=2\pi k_e\lambda R
\frac{x}{(R^2+x^2)^{3/2}}.
\]

为求极值，对 \(x(R^2+x^2)^{-3/2}\) 求导：

\[
\begin{aligned}
\frac{d}{dx}\left[x(R^2+x^2)^{-3/2}\right]
&=(R^2+x^2)^{-3/2}\\
&\quad-3x^2(R^2+x^2)^{-5/2}\\
&=\frac{R^2-2x^2}{(R^2+x^2)^{5/2}}.
\end{aligned}
\]

极值条件为

\[
R^2-2x^2=0.
\]

在 \(x>0\) 取

\[
\boxed{x_{\max}=\frac{R}{\sqrt2}}.
\]

由于 \(E(0)=0\) 且 \(x\to\infty\) 时 \(E\to0\)，这个驻点是最大值。代入：

\[
\begin{aligned}
E_{\max}
&=2\pi k_e\lambda R
\frac{R/\sqrt2}{(3R^2/2)^{3/2}}\\
&=\frac{4\sqrt3\pi k_e\lambda}{9R}\\
&=\frac{\lambda}{3\sqrt3\,\varepsilon_0R}.
\end{aligned}
\]

所以

\[
\boxed{E_{\max}=\frac{\lambda}{3\sqrt3\,\varepsilon_0R}}.
\]

若 \(\lambda>0\)，方向沿 \(+x\)；若
\(\lambda<0\)，方向相反，上式给出大小。

\subsubsection{2.
半球弧面在原点产生的电场}\label{ux534aux7403ux5f27ux9762ux5728ux539fux70b9ux4ea7ux751fux7684ux7535ux573a}

把图 6-2(b) 中的场点 \(P\) 移到原点 \(O\)，便得到第 2
问的几何关系。取半球对称轴的负方向为极轴，令 \(\theta\) 表示面元半径与负
\(x\) 轴的夹角，范围为

\[
0\le\theta\le\frac{\pi}{2}.
\]

半球上极角在 \(\theta\) 到 \(\theta+d\theta\) 之间的窄带，宽度为
\(R\,d\theta\)，纬圈半径为 \(R\sin\theta\)，周长为
\(2\pi R\sin\theta\)。因此图 6-2(b) 中的面元带面积为

\[
\begin{aligned}
dS
&=(\text{纬圈周长})\times(\text{窄带宽度})\\
&=(2\pi R\sin\theta)(R\,d\theta)\\
&=2\pi R^2\sin\theta\,d\theta.
\end{aligned}
\]

原点到任一面元距离都是 \(R\)，面元产生的场强大小为

\[
dE=k_e\frac{\sigma\,dS}{R^2}.
\]

同一纬圈上的面元关于 \(x\)
轴对称，所有横向分量相消，轴向分量相加。由图中的投影关系，轴向分量为

\[
dE_x=dE\cos\theta.
\]

于是

\[
\begin{aligned}
E_O
&=\int_0^{\pi/2}k_e\frac{\sigma}{R^2}
\left(2\pi R^2\sin\theta\,d\theta\right)\cos\theta\\
&=2\pi k_e\sigma\int_0^{\pi/2}\sin\theta\cos\theta\,d\theta\\
&=2\pi k_e\sigma\cdot\frac12\\
&=\pi k_e\sigma.
\end{aligned}
\]

因此

\[
\boxed{
\bm E_O=\frac{\sigma}{4\varepsilon_0}\,\hat{\bm x}
}.
\]

\subsubsection{3.
半球弧面在轴上任意点的场}\label{ux534aux7403ux5f27ux9762ux5728ux8f74ux4e0aux4efbux610fux70b9ux7684ux573a}

现在把图 6-2(b) 中的场点放在 \(P(x_0,0,0)\)，记 \(x=x_0\ge0\)。仍以负
\(x\) 轴为极轴，沿用第 2 问的极角 \(\theta\)，所以

\[
0\le\theta\le\frac{\pi}{2}.
\]

图中半球位于原点左侧。极角为 \(\theta\)
的面元，其轴坐标和到轴线的距离分别为

\[
x_S=-R\cos\theta,
\qquad
\rho_S=R\sin\theta.
\]

所以，从面元 \(S\) 指向场点 \(P\) 的矢量，其轴向分量为

\[
x-x_S=x+R\cos\theta.
\]

图 6-2(b) 中的距离 \(s\) 由轴向距离与横向距离构成直角三角形，故

\[
\begin{aligned}
s^2
&=(x+R\cos\theta)^2+(R\sin\theta)^2\\
&=x^2+2xR\cos\theta+R^2\cos^2\theta+R^2\sin^2\theta\\
&=x^2+R^2+2xR\cos\theta,
\end{aligned}
\]

即

\[
s=\sqrt{x^2+R^2+2xR\cos\theta}.
\]

一个面元带所带电量为 \(dq=\sigma\,dS\)，其场强大小为
\(dE=k_e\sigma\,dS/s^2\)。由图中的分量三角形，轴向投影因子为
\((x+R\cos\theta)/s\)，因此

\[
\begin{aligned}
dE_x
&=dE\frac{x+R\cos\theta}{s}\\
&=k_e\frac{\sigma\,dS}{s^2}
\frac{x+R\cos\theta}{s}\\
&=k_e\sigma\frac{x+R\cos\theta}{s^3}
\left(2\pi R^2\sin\theta\,d\theta\right).
\end{aligned}
\]

至此，图中的三个几何量已经逐一进入公式：\(dS\) 决定电荷量，\(s\)
决定平方反比因子，\(x+R\cos\theta\) 决定轴向投影。对
\(0\le\theta\le\pi/2\) 的全部纬圈求和，得到

\[
E_x=2\pi k_e\sigma R^2
\int_0^{\pi/2}
\frac{(x+R\cos\theta)\sin\theta\,d\theta}
{(x^2+R^2+2xR\cos\theta)^{3/2}}.
\]

令

\[
u=\cos\theta,
\qquad
du=-\sin\theta\,d\theta.
\]

当 \(\theta=0\) 时 \(u=1\)，当 \(\theta=\pi/2\) 时
\(u=0\)。负号与上下限交换相互抵消，于是

\[
E_x=2\pi k_e\sigma R^2
\int_0^1
\frac{x+Ru}{(x^2+R^2+2xRu)^{3/2}}\,du.
\]

再引入无量纲变量

\[
a=\frac{x}{R}.
\]

约去相同的 \(R\) 次幂：

\[
E_x=2\pi k_e\sigma
\int_0^1
\frac{a+u}{(a^2+1+2au)^{3/2}}\,du.
\]

以下换元先按 \(a>0\) 进行；\(a=0\) 已由第 2
问单独求得，最后也可由连续极限恢复。令

\[
w=a^2+1+2au.
\]

于是

\[
du=\frac{dw}{2a},
\qquad
u=\frac{w-a^2-1}{2a}.
\]

分子可写成

\[
a+u
=a+\frac{w-a^2-1}{2a}
=\frac{a^2-1+w}{2a}.
\]

所以不定积分化为

\[
\begin{aligned}
\int\frac{a+u}{(a^2+1+2au)^{3/2}}\,du
&=\frac{1}{4a^2}\int
\left[(a^2-1)w^{-3/2}+w^{-1/2}\right]dw\\
&=-\frac{a^2-1}{2a^2\sqrt w}
+\frac{\sqrt w}{2a^2}+C.
\end{aligned}
\]

积分端点 \(u=0,1\) 分别对应

\[
w_0=a^2+1,
\qquad
w_1=(a+1)^2.
\]

因为 \(a>0\)，所以 \(\sqrt{w_1}=a+1\)。记上述原函数为 \(F(w)\)，在上限处

\[
\begin{aligned}
F(w_1)
&=-\frac{a^2-1}{2a^2(a+1)}
+\frac{a+1}{2a^2}\\
&=-\frac{a-1}{2a^2}+\frac{a+1}{2a^2}\\
&=\frac{1}{a^2}.
\end{aligned}
\]

在下限处

\[
\begin{aligned}
F(w_0)
&=-\frac{a^2-1}{2a^2\sqrt{1+a^2}}
+\frac{\sqrt{1+a^2}}{2a^2}\\
&=\frac{-(a^2-1)+(1+a^2)}
{2a^2\sqrt{1+a^2}}\\
&=\frac{1}{a^2\sqrt{1+a^2}}.
\end{aligned}
\]

因此定积分为

\[
\begin{aligned}
F(w_1)-F(w_0)
&=\frac{1}{a^2}
-\frac{1}{a^2\sqrt{1+a^2}}\\
&=\frac{\sqrt{1+a^2}-1}
{a^2\sqrt{1+a^2}}.
\end{aligned}
\]

再利用有理化恒等式

\[
\sqrt{1+a^2}-1
=\frac{a^2}{\sqrt{1+a^2}+1},
\]

便得到前述紧凑结果。这里把代数变形完整展开，是为了明确每一项的来源并便于检查符号。

所以

\[
E_x
=\frac{2\pi k_e\sigma}
{\sqrt{1+a^2}\left(1+\sqrt{1+a^2}\right)}.
\]

还原 \(a=x/R\)，可写成

\[
\boxed{
\bm E(x)=
\frac{2\pi k_e\sigma R^2}
{R^2+x^2+R\sqrt{R^2+x^2}}
\,\hat{\bm x}
}.
\]

检验 \(x=0\)：

\[
E(0)=\frac{2\pi k_e\sigma R^2}{2R^2}=\pi k_e\sigma,
\]

与第2问一致。远处 \(x\gg R\) 时，分母主项为 \(x^2\)，于是

\[
E(x)\sim\frac{2\pi k_e\sigma R^2}{x^2}
=k_e\frac{Q}{x^2},
\]

其中 \(Q=2\pi R^2\sigma\) 正是半球弧面的总电荷，远场检验通过。

还应检查方向。对 \(x\ge0\)，每个正电面元指向场点的矢量都有正的 \(x\)
分量，所以上式分母始终为正且
\(E_x>0\)。若轴上结果发生方向翻转，就说明极角范围或投影符号设置错误。

\newpage

\section{第 7 题}\label{ux7b2c-7-ux9898}

一个半径为 \(R\)
的光滑细圆环固定在竖直平面内。一根光滑、柔软、不可伸长的闭合细链挂在圆环上，链中串有一个质量为
\(m\) 的光滑小物体。链的两个接触点为 \(A\)、\(B\)，小物体位于
\(C\)，装置关于过圆心 \(O\) 的竖直线对称。记链段 \(AB\) 与半径 \(AO\)
的夹角为 \(\theta\)。

\begin{enumerate}
\def\labelenumi{\arabic{enumi}.}
\tightlist
\item
  求闭合链能够挂在圆环上的最小总长度 \(L_0\)；
\item
  平衡时若 \(\theta=30^\circ\)，求链的总长度 \(L_1\)；
\item
  接第1问，若小物体在圆环所在平面内作小振动，求周期 \(T_1\)。
\end{enumerate}

\begin{figure}[H]
\centering
\includegraphics[width=0.493\linewidth,height=\textheight,keepaspectratio,alt={图 7-1　闭合柔链在竖直圆环上的对称悬挂状态。}]{../figures/question-07.pdf}
\caption*{图 7-1　闭合柔链在竖直圆环上的对称悬挂状态。}
\end{figure}

\subsection{一、题型定位}\label{ux4e00ux9898ux578bux5b9aux4f4d-6}

轻链的等张力与接触力方向共同限制几何构型。先确定可行链长，再求下端的等效摆长；退化构型的无穷张力和第3问缺少的角度条件都要保留。

\subsection{二、知识点}\label{ux4e8cux77e5ux8bc6ux70b9-6}

\begin{itemize}
\tightlist
\item
  光滑无质量链的等张力；
\item
  光滑圆环的径向支持力；
\item
  受力角平分线与长度约束；
\item
  退化构型和有限张力；
\item
  固定长度下的圆周轨迹；
\item
  小振动的动能、势能与有效摆长。
\end{itemize}

\subsection{三、物理图像}\label{ux4e09ux7269ux7406ux56feux50cf-6}

\(A,B\)
关于竖直轴对称，位于圆心上方。两等张力的合力指向圆心，圆环对链的支持力反向、沿半径向外。链段
\(AC\)、\(BC\) 与水平的夹角为 \(2\theta\)。物体 \(C\) 位于圆心下方，距离
\(OC\) 由平衡角决定。

链长固定以后，三角形 \(ABC\)
的形状固定；接触点可沿光滑环移动，整个三角形在圆环平面内转动。质量点的运动轨迹是以
\(O\) 为圆心、\(OC\) 为半径的圆，而不是以链上另一个点为支点的摆。

\subsection{四、解题入口}\label{ux56dbux89e3ux9898ux5165ux53e3-6}

先由水平投影求 \(AB,AC\)，组成 \(L(\theta)\) 并证明它单调。再由 \(A\)
的高度减去 \(AC\) 的竖直投影求 \(C\) 的位置。固定非退化链长对应固定
\(\theta\)，以绕 \(O\) 的小转角 \(\varphi\) 建立振动方程。

\subsection{五、解析}\label{ux4e94ux89e3ux6790-6}

\begin{figure}[H]
\centering
\includegraphics[width=0.771\linewidth,height=\textheight,keepaspectratio,alt={图 7-2　接触点三力、固定链长下的圆周轨迹及小振动。}]{../figures/solution-07.pdf}
\caption*{图 7-2　接触点三力、固定链长下的圆周轨迹及小振动。}
\end{figure}

图 7-2 左图中的 \(N\) 与两张力合力方向相反；右图中支点就是圆心
\(O\)，有效摆长为 \(\ell=OC\)。以下计算分别从这两幅图取得角度和长度。

\textbf{1. 链的长度}

由 \(OA=OB=R\)，半弦长为
\(R\cos\theta\)。接触点两张力等大，半径平分它们的夹角，所以

\[AB=2R\cos\theta,\qquad AC=BC=\frac{R\cos\theta}{\cos2\theta}.\]

从而

\[\boxed{L(\theta)=2R\cos\theta(1+\sec2\theta)=\frac{4R\cos^3\theta}{\cos2\theta}},\quad0<\theta<\pi/4.\]

正的角度保证有向下承重的链段；\(\cos2\theta>0\)
保证段长为正。对长度求导：

\[L'(\theta)=\frac{4R(3-2\cos^2\theta)\sin\theta\cos^2\theta}{(2\cos^2\theta-1)^2}>0.\]

因此长度的下确界为

\[\boxed{L_0=4R}.\]

这里 \(\theta\to0^+\) 时三点趋于共线。物体的平衡条件
\(2T\sin2\theta=mg\) 给 \(T\to\infty\)，所以有限张力的真实承重构型必须有
\(L>4R\)。原参考中的``最小长度''是这一退化极限，不能将它当作普通平衡点。

\textbf{2. \(30^\circ\) 平衡角的长度}

代入 \(\cos\theta=\sqrt3/2\)、\(\cos2\theta=1/2\)：

\[\boxed{L_1=3\sqrt3R}.\]

此时 \(AB=\sqrt3R\)，另外两边也各为 \(\sqrt3R\)，三角形为等边三角形。

\textbf{3. 一般平衡角的有效摆长}

取圆心为原点、竖直向上为正，\(A\) 高度为 \(R\sin\theta\)。物体坐标为

\[y_C=R\sin\theta-AC\sin2\theta=-\frac{R\sin\theta}{\cos2\theta}.\]

所以

\[\boxed{\ell=OC=\frac{R\sin\theta}{\cos2\theta}}.\]

由 \(L'(\theta)>0\)，每个固定的 \(L>4R\) 唯一确定这个非退化角。整体转过
\(\varphi\) 后，\(C\) 位于同一半径的圆上。

\textbf{4. 小振动方程与周期}

链与圆环均无质量且光滑，只有物体贡献动能，圆环约束不做功：

\[K=\frac12m\ell^2\dot\varphi^2,\qquad U=mg\ell(1-\cos\varphi).\]

拉格朗日方程为 \(m\ell^2\ddot\varphi+mg\ell\sin\varphi=0\)。取
\(|\varphi|\ll1\) 得

\[\boxed{T_1=2\pi\sqrt{\frac{R\sin\theta}{g\cos2\theta}}}.\]

若第3问意在承接第2问的 \(\theta=30^\circ\)，才有 \(\ell=R\)，于是

\[\boxed{T_1=2\pi\sqrt{R/g}}\qquad(\theta=30^\circ).\]

只承接第1问的长度下确界，并不能得到这一周期；\(\theta=0\)
处张力发散、构型退化，普通单摆线性化不适用。

\textbf{5. 极限核对}

小角展开 \(L=4R+2R\theta^2+O(\theta^4)\)，而
\(\ell\sim R\theta\)。因此在有限非零小角处，\(T\sim2\pi\sqrt{R\theta/g}\)，趋近退化边界时变短；这与固定摆长
\(R\) 的错误推断不同。另一方面 \(\theta\to\pi/4^-\)
时链长和有效摆长都发散，周期也发散。周期与物体质量无关，符合单自由度重力摆的量纲检查。

\newpage

\section{第 8 题}\label{ux7b2c-8-ux9898}

空气净化器由筒状滤芯和上部离心风扇组成。滤芯内径为
\(r_2\)，底部密封。风扇由 \(9\) 片扇叶组成，每片扇叶的内、外半径分别为
\(r_1\)、\(r_2\)，可视为与水平面夹角为 \(\theta\)、宽度为 \(d\)
的平面，工作时绕中心 \(O\) 以角速度 \(\omega\) 转动。

采用如下理想模型：滤芯桶内空气始终均匀，温度恒为
\(T\)；碰撞前空气分子相对实验室静止，扇叶扫过的空气全部被排到外界；滤芯单位时间的通气摩尔数与内外压强差成正比，

\[
\frac{\Delta n}{\Delta t}=k\Delta p.
\]

外界大气密度为 \(\rho_0\)、压强为 \(p_0\)，内外空气均视为理想气体。

\begin{enumerate}
\def\labelenumi{\arabic{enumi}.}
\tightlist
\item
  风扇角速度为 \(\omega\) 时，求稳定筒内压强比 \(p/p_0\)；
\item
  求每秒通风体积 \(V\)，体积折算到外界温度和大气压强；
\item
  求风机对扇叶的总输出功率 \(P\)。
\end{enumerate}

数值取

\[
\begin{gathered}
k=0.002\,\mathrm{mol/(s\cdot Pa)},\quad
R_g=8.31\,\mathrm{J/(mol\cdot K)},\quad T=300\,\mathrm K,\\
\omega=70\,\mathrm{rad/s},\quad d=0.02\,\mathrm m,
\quad \rho_0=1.29\,\mathrm{kg/m^3},\\
\theta=0.9\,\mathrm{rad},\quad r_1=0.07\,\mathrm m,
\quad r_2=0.20\,\mathrm m.
\end{gathered}
\]

\begin{figure}[H]
\centering
\includegraphics[width=0.783\linewidth,height=\textheight,keepaspectratio,alt={图 8-1　筒状滤芯、九叶离心风扇及扇叶倾角。}]{../figures/question-08.pdf}
\caption*{图 8-1　筒状滤芯、九叶离心风扇及扇叶倾角。}
\end{figure}

\subsection{一、题型定位}\label{ux4e00ux9898ux578bux5b9aux4f4d-7}

先算滤芯的体积流率，再算静止空气与扇叶碰撞后获得的速度和动量。用进出流量关系确定压力及排风量，并把动能增长与机械功率区分开。

\subsection{二、知识点}\label{ux4e8cux77e5ux8bc6ux70b9-7}

\begin{itemize}
\tightlist
\item
  扇叶扫掠体积；
\item
  稳态摩尔数守恒；
\item
  理想气体摩尔密度 \(p/(R_gT)\)；
\item
  不同压强下体积流量的换算；
\item
  移动斜壁的弹性碰撞；
\item
  功率的微元积分。
\end{itemize}

\subsection{三、物理图像}\label{ux4e09ux7269ux7406ux56feux50cf-7}

风扇把筒内空气向外排出，使筒内压强 \(p\) 低于外界
\(p_0\)。滤芯通过压强差吸入空气，稳态时``滤芯流入摩尔率''恰好等于``风扇排出摩尔率''。扇叶在半径
\(r\) 处的线速度为 \(\omega r\)，其垂直于扇叶的分量为
\(\omega r\sin\theta\)。

\subsection{四、解题入口}\label{ux56dbux89e3ux9898ux5165ux53e3-7}

先把九片扇叶每秒扫出的几何体积写成
\(\dot V_f\)。排出的摩尔流率不是简单的 \(\dot V_f\)，而是筒内摩尔密度
\(p/(R_gT)\) 与 \(\dot V_f\) 的乘积。第三问先求半径 \(r\) 处、宽度
\(dr\) 的扇叶条带在 \(dt\) 内碰到的空气质量，再乘每千克空气获得的动能。

\subsection{五、解析}\label{ux4e94ux89e3ux6790-7}

\begin{figure}[H]
\centering
\includegraphics[width=0.821\linewidth,height=\textheight,keepaspectratio,alt={图 8-2　稳态摩尔流率平衡与移动扇叶对空气的碰撞。}]{../figures/solution-08.pdf}
\caption*{图 8-2　稳态摩尔流率平衡与移动扇叶对空气的碰撞。}
\end{figure}

图 8-2(a) 区分了两种摩尔流率：滤芯流入为 \(k(p_0-p)\)，风扇排出为
\(p\dot V_f/(R_gT)\)。图 8-2(b) 表明扇叶线速度为
\(\omega r\)，其法向分量为
\(v_n=\omega r\sin\theta\)；移动弹性壁使静止空气获得的法向速度为
\(2v_n\)。后面的质量微元和动能微元都由这两个图中量建立。

\subsubsection{1.
风扇的几何排风体积率}\label{ux98ceux6247ux7684ux51e0ux4f55ux6392ux98ceux4f53ux79efux7387}

在时间 \(dt\) 内，一片扇叶转过角度

\[
d\phi=\omega\,dt.
\]

从 \(r_1\) 到 \(r_2\) 的环形区域面积为

\[
A_{12}=\frac12(r_2^2-r_1^2)d\phi.
\]

扇叶宽度在排风方向上的有效投影为
\(d\sin\theta\)。因此一片扇叶扫出的有效体积为

\[
dV_1=\frac12(r_2^2-r_1^2)\omega\,dt\,d\sin\theta.
\]

共有 \(9\) 片扇叶，所以

\[
dV_f=\frac92\omega d(r_2^2-r_1^2)\sin\theta\,dt.
\]

定义

\[
A=9\omega d(r_2^2-r_1^2)\sin\theta,
\]

则风扇几何排风体积率为

\[
\boxed{\dot V_f=\frac{A}{2}}.
\]

\subsubsection{2.
稳态筒内压强}\label{ux7a33ux6001ux7b52ux5185ux538bux5f3a}

滤芯在单位时间内流入筒内的摩尔数为

\[
\dot n_{\rm in}=k(p_0-p).
\]

筒内空气的摩尔密度为

\[
c_n=\frac{p}{R_gT}.
\]

风扇排出的摩尔流率为

\[
\dot n_{\rm out}=c_n\dot V_f
=\frac{p}{R_gT}\frac{A}{2}.
\]

稳态条件是

\[
\dot n_{\rm in}=\dot n_{\rm out}.
\]

所以

\[
k(p_0-p)=\frac{pA}{2R_gT}.
\]

两边乘以 \(2R_gT\)：

\[
2kR_gTp_0-2kR_gTp=pA.
\]

把含 \(p\) 的项移到右侧：

\[
2kR_gTp_0=p(2kR_gT+A).
\]

因此

\[
\boxed{
\frac{p}{p_0}
=\frac{2kR_gT}
{2kR_gT+9\omega d(r_2^2-r_1^2)\sin\theta}
}.
\]

代入数值：

\[
A=9\times70\times0.02\times(0.20^2-0.07^2)\sin0.9
\approx0.34643\,\mathrm{m^3/s}.
\]

\[
2kR_gT=2\times0.002\times8.31\times300
\approx9.972\,\mathrm{m^3/s}.
\]

所以

\[
\boxed{\frac{p}{p_0}\approx0.96643}.
\]

筒内只需维持约 \(3.36\%\) 的压强降低，即可驱动滤芯进气。

\subsubsection{3.
折算到外界状态的通风体积}\label{ux6298ux7b97ux5230ux5916ux754cux72b6ux6001ux7684ux901aux98ceux4f53ux79ef}

风扇实际从筒内排出的几何体积率为
\(\dot V_f=A/2\)，但题目要求把体积折算到外界压强 \(p_0\)、同一温度
\(T\)。相同摩尔流率在外界状态下的体积率为

\[
V=\dot n_{\rm out}\frac{R_gT}{p_0}.
\]

又因为

\[
\dot n_{\rm out}=\frac{p}{R_gT}\dot V_f,
\]

所以

\[
V=\frac{p}{p_0}\dot V_f.
\]

代入 \(p/p_0\) 与 \(\dot V_f=A/2\)：

\[
\boxed{
V=\frac{kR_gT\,A}{2kR_gT+A}
}.
\]

展开 \(A\)：

\[
\boxed{
V=
\frac{9kR_gT\,\omega d(r_2^2-r_1^2)\sin\theta}
{2kR_gT+9\omega d(r_2^2-r_1^2)\sin\theta}
}.
\]

数值为

\[
\dot V_f=\frac{A}{2}\approx0.17322\,\mathrm{m^3/s},
\]

\[
\boxed{V\approx0.16740\,\mathrm{m^3/s}}.
\]

\subsubsection{4.
筒内空气密度}\label{ux7b52ux5185ux7a7aux6c14ux5bc6ux5ea6}

内外温度相同，理想气体密度与压强成正比，因此

\[
\rho=\rho_0\frac{p}{p_0}.
\]

即

\[
\rho=\rho_0\frac{2kR_gT}{2kR_gT+A}.
\]

这个密度必须用于功率积分；若直接使用外界密度
\(\rho_0\)，会高估扇叶碰到的空气质量。

\subsubsection{5.
单个扇叶微元传给空气的动能}\label{ux5355ux4e2aux6247ux53f6ux5faeux5143ux4f20ux7ed9ux7a7aux6c14ux7684ux52a8ux80fd}

取一片扇叶上半径为 \(r\)、径向宽度为 \(dr\) 的条带。条带面积为

\[
dS=d\,dr.
\]

扇叶速度为 \(\omega r\)，垂直于扇叶的速度分量为

\[
v_n=\omega r\sin\theta.
\]

在时间 \(dt\) 内，与该条带碰撞的空气体积为

\[
d\mathcal V=v_n\,dt\,dS
=\omega r\sin\theta\,dt\,d\,dr.
\]

对应空气质量为

\[
dm=\rho\,d\mathcal V
=\rho\omega r\sin\theta\,d\,dr\,dt.
\]

按照题设的弹性碰撞模型，静止空气分子与以法向速度 \(v_n\)
运动的扇叶碰撞后，法向速度变为

\[
v_{\rm air}=2v_n=2\omega r\sin\theta.
\]

因此该空气微元获得的动能为

\[
\begin{aligned}
dW_1
&=\frac12dm\,v_{\rm air}^2\\
&=\frac12\left(\rho\omega r\sin\theta\,d\,dr\,dt\right)
\left(2\omega r\sin\theta\right)^2\\
&=2\rho d\omega^3r^3\sin^3\theta\,dr\,dt.
\end{aligned}
\]

\subsubsection{6.
九片扇叶的总功率}\label{ux4e5dux7247ux6247ux53f6ux7684ux603bux529fux7387}

对 \(r_1\) 到 \(r_2\) 积分，并乘以 \(9\)：

\[
\begin{aligned}
dW
&=9\int_{r_1}^{r_2}
2\rho d\omega^3r^3\sin^3\theta\,dr\,dt\\
&=18\rho d\omega^3\sin^3\theta
\left[\frac{r^4}{4}\right]_{r_1}^{r_2}dt\\
&=\frac92\rho d\omega^3\sin^3\theta
(r_2^4-r_1^4)dt.
\end{aligned}
\]

所以功率为

\[
P=\frac{dW}{dt}
=\frac92\rho d\omega^3\sin^3\theta(r_2^4-r_1^4).
\]

代入 \(\rho=\rho_0p/p_0\)：

\[
\boxed{
P=
\frac{9kR_gT\,\rho_0d\omega^3\sin^3\theta(r_2^4-r_1^4)}
{2kR_gT+9\omega d(r_2^2-r_1^2)\sin\theta}
}.
\]

代入数值：

\[
\boxed{P\approx29.15\,\mathrm W}.
\]

量纲检查：\(\rho d\omega^3r^4\) 的量纲为

\[
\frac{\mathrm{kg}}{\mathrm{m^3}}\cdot\mathrm m
\cdot\frac{1}{\mathrm{s^3}}\cdot\mathrm{m^4}
=\frac{\mathrm{kg\,m^2}}{\mathrm{s^3}},
\]

正是功率单位瓦特。

\subsubsection{7.
极限与数量级复核}\label{ux6781ux9650ux4e0eux6570ux91cfux7ea7ux590dux6838}

若滤芯极易通气，即 \(k\to\infty\)，则

\[
\frac{p}{p_0}\to1.
\]

此时筒内几乎不形成压降，折算通风体积趋于风扇的几何排风量：

\[
V\to\frac{A}{2}.
\]

若滤芯完全堵塞，即 \(k\to0\)，则

\[
\frac{p}{p_0}\to0,
\qquad
V\to0.
\]

在理想模型中风扇会持续抽走有限筒内空气，使稳态压强趋近于零。真实装置会受到漏气、回流和风扇性能曲线限制，因此不能把这一极限解释为实际真空能力。

若 \(\omega\to0\)，则 \(A\to0\)，压强比趋于 \(1\)、通风量趋于
\(0\)、功率按 \(\omega^3\) 趋于 \(0\)，三项结论相互一致。

数值上，几何排风率和折算通风量分别为

\[
\dot V_f\approx0.17322\,\mathrm{m^3/s},
\qquad
V\approx0.16740\,\mathrm{m^3/s}.
\]

二者之比

\[
\frac{V}{\dot V_f}
\approx\frac{0.16740}{0.17322}
\approx0.9664
\]

恰好等于 \(p/p_0\)，这是第1、2问之间最直接的数值交叉检查。
\end{document}
